% !TeX program = xelatex
% UTF-8. Compile with XeLaTeX twice (or use latexmk -xelatex).
% Self-contained source: no external figures, bibliography or input files.
\IfFileExists{ctexart.cls}{
  \documentclass[UTF8,11pt,a4paper]{ctexart}
}{
  \documentclass[11pt,a4paper]{article}
  \usepackage{fontspec}
  \IfFontExistsTF{Noto Sans CJK SC}{
    \setmainfont[AutoFakeBold=2,AutoFakeSlant=0.15]{Noto Sans CJK SC}
  }{
    \setmainfont[AutoFakeBold=2,AutoFakeSlant=0.15]{Microsoft YaHei}
  }
  \XeTeXlinebreaklocale "zh"
  \XeTeXlinebreakskip=0pt plus 1pt
}
\usepackage[margin=23mm,headheight=16pt,headsep=7mm,footskip=12mm]{geometry}
\usepackage{amsmath,amssymb,amsthm,mathtools,mathrsfs}
\usepackage{booktabs,array,longtable,fancyhdr,enumitem}
\usepackage[dvipsnames]{xcolor}
\usepackage[hidelinks,bookmarksnumbered,bookmarksdepth=2,hypertexnames=false]{hyperref}
\hypersetup{pdftitle={紧线性差与自然函数方程：详细汇报},pdfauthor={XP-算子项目},pdfsubject={存在唯一性、自然性、有限模板及四点支撑分类}}
\theoremstyle{definition}
\newtheorem{theorem}{定理}[section]
\newtheorem{proposition}[theorem]{命题}
\newtheorem{lemma}[theorem]{引理}
\newtheorem{corollary}[theorem]{推论}
\newtheorem{definition}[theorem]{定义}
\newtheorem{example}[theorem]{例}
\newtheorem{remark}[theorem]{说明}
\renewcommand{\proofname}{\normalfont\bfseries 证明}
\renewcommand{\contentsname}{目录}
\renewcommand{\refname}{参考文献}
\newcommand{\K}{\mathbb K}
\newcommand{\Q}{\mathbb Q}
\newcommand{\R}{\mathbb R}
\newcommand{\C}{\mathbb C}
\newcommand{\Z}{\mathbb Z}
\newcommand{\N}{\mathbb N}
\newcommand{\Hom}{\operatorname{Hom}}
\newcommand{\End}{\operatorname{End}}
\newcommand{\im}{\operatorname{im}}
\newcommand{\coker}{\operatorname{coker}}
\newcommand{\Span}{\operatorname{span}}
\newcommand{\supp}{\operatorname{supp}}
\newcommand{\id}{\operatorname{id}}
\newcommand{\ev}{\operatorname{ev}}
\newcommand{\Ann}{\operatorname{Ann}}
\newcommand{\D}{\mathcal D}
\newcommand{\T}{\mathcal T}
\newcommand{\calA}{\mathcal A}
\newcommand{\calM}{\mathcal M}
\newcommand{\calN}{\mathcal N}
\newcommand{\Sol}{\operatorname{Sol}}
\newcommand{\Uc}{\mathscr U_c}
\newcommand{\Wc}{\mathscr W_c}
\DeclareMathOperator{\rank}{rank}
\setlength{\parskip}{0.3em}
\setlength{\emergencystretch}{3em}
\renewcommand{\arraystretch}{1.2}
\setlist[enumerate]{itemsep=2pt,topsep=4pt}
\setlist[itemize]{itemsep=2pt,topsep=4pt}
\allowdisplaybreaks[1]
\numberwithin{equation}{section}
\pagestyle{fancy}
\fancyhf{}
\fancyhead[L]{\small XP-算子：紧线性差与自然函数方程}
\fancyhead[R]{\small 详细汇报}
\fancyfoot[C]{\thepage}
\begin{document}
\begin{titlepage}
\vspace*{16mm}
{\small\color{MidnightBlue} XP-算子项目\par}
\vspace{8mm}
{\LARGE\bfseries 紧线性差与自然函数方程\par}
\vspace{5mm}
{\Large 详细汇报\par}
\vspace{8mm}
{\large 存在唯一性、有限模板与低支撑分类\par}
\vspace{10mm}
{\normalsize 2026年10月9日\par}
\vspace{10mm}
\noindent\textbf{摘要。}
本报告统一采用修订后的术语：线性差是对函数变量的 $K$-线性算子，仅满足可加性的对象称为加性差。
首先证明紧线性差的局部线性表示，给出 $Q=P-h$ 全域函数问题的存在性、唯一性和全部解的参数化；明确标量值、有限维向量值与一般无限维值域的不同结论。
其次区分固定空间的群等变性、自同态相容性与跨空间的完全线性自然性，并由后者得到统一的有限取值模板。
在此基础上，整理模板矩、一元伸缩、乘法群代数与有限域的计算模型。
报告的重点是特征零、模板秩二、总支撑至多四点的全部齐次解分类：单点方向给出广义多项式次数上界；三个、四个方向分别归结为加性映射和对称双加性映射；两个方向各有两点时，由乘法轨道和字符控制全部解。
最后用两个交换代数商将结构完全参数化，给出代数参数的有限矩阵公式，以及“恒等映射加形式导子”和“共轭嵌入加范数”的两个完整算例。

\medskip
\noindent\textbf{阅读约定。}
函数均为任意集合映射，除特别说明外不附加连续性、可测性或普通多项式假设。
星号表示完整的 $K$-线性代数对偶，而不是连续对偶或环同态集合。
使用通常的选择公理以选取基、补空间和线性延拓；$\N=\{0,1,2,\ldots\}$。
“紧”采用本项目的有限取值可分解定义，不指泛函分析中的紧算子。

\medskip
\noindent\textbf{整理范围。}
报告围绕本会话的线性差、自然性与低支撑研究组织，补齐论证所需的基础命题。
各定理均在正文给出证明；零空间、特征、维数及参数域的限制写入相应结论。
非齐次存在性的相容判据与齐次核分类分别陈述。
\vfill
\noindent\small 编译：单文件 UTF-8 源码；使用 XeLaTeX 连续编译两次即可生成目录、书签及交叉引用。
\end{titlepage}
\pagenumbering{roman}
\setcounter{tocdepth}{1}
\begingroup
\normalsize
\setlength{\parskip}{0pt}
\tableofcontents
\endgroup
\clearpage
\pagenumbering{arabic}
\section{基本对象、问题类型与线性性的层次}
设 $R,S$ 为交换环，$V,W$ 分别为 $R,S$-模，$n\ge1$，并令
\[
 I=V,\qquad J=V^n,\qquad W^I=\{f:I\to W\}.
\]
对映射 $P:W^I\times J\to W$，定义
\[
 X_P(f)=\{v\in J:P(f,v)=0\}.
\]
它称为 $f$ 的 $P$-域，也沿用讨论中的 $P$-Area 一词；这里的“域”不指代数意义的域。

\begin{definition}[三类问题与本稿的非齐次问题]
$P$-函数问题要求 $X_P(f)=J$；反 $P$-函数问题要求 $X_P(f)=\varnothing$；
$P$-域问题要求对给定 $f$ 求 $X_P(f)$。
给定 $h:J\to W$，令 $Q(f,v)=P(f,v)-h(v)$。本稿所说的 $Q$-问题均指
\begin{equation}\label{eq:inhom}
 P(f,v)=h(v)\qquad(v\in J),
\end{equation}
即全域 $Q$-函数问题。要求 $X_Q(f)$ 恰为某个指定真子集，是另一个实现问题。
\end{definition}

\begin{definition}[加性差与线性差：修订后的约定]
若对全部 $f,g,v$ 有
\[
 P(f+g,v)=P(f,v)+P(g,v),
\]
则称 $P$ 为\textbf{加性差}。
当 $R=S=K$ 为域时，若此外对每个 $a\in K$ 还满足
\[
 P(af,v)=aP(f,v),
\]
则称 $P$ 为\textbf{线性差}，必要时强调为 $K$-线性差。
本文统一采用这一修订，不再用“线性差”指仅满足可加性的对象。
\end{definition}
可加性使 $P(0,v)=0$。特征零时，可加映射自动为 $\Q$-线性；
特征 $p$ 时自动为 $\mathbb F_p$-线性。一般并不自动为 $K$-线性。
例如 $K=\C$ 时 $P(f,v)=\overline{f(0)}$ 可加，却不是复线性的。

\begin{proposition}[一般非齐次方程的纤维]
若 $P$ 为 $K$-线性差，令
\[
 \T:W^I\to W^J,\qquad (\T f)(v)=P(f,v).
\]
则方程\eqref{eq:inhom}可解当且仅当 $h\in\im\T$；有一个解 $f_0$ 时，全部解为
\begin{equation}\label{eq:torsor}
 \Sol(h)=f_0+\ker\T.
\end{equation}
对某个可解的 $h$ 解唯一，当且仅当 $\ker\T=0$；这也等价于每个可解的右端均有唯一解。
每个右端恰有一解，当且仅当 $\T$ 为同构。
\end{proposition}
\begin{proof}
$\T f=h=\T f_0$ 当且仅当 $\T(f-f_0)=0$，得到\eqref{eq:torsor}。
存在性、唯一性和全体右端的结论分别就是像、核及双射的定义。
\end{proof}
若只假设可加性，同一证明在加法群层次成立，在素域层次给出仿射空间。
即便 $V,W$ 有限维，$W^V$ 通常仍是巨大的函数空间；不能对 $V,W$ 的有限维数直接套用一个有限矩阵的秩公式。

\section{紧算子的完整等价刻画}\label{sec:compact}
本节只需要 $W$ 是阿贝尔群，$I\ne\varnothing$。对 $x\in I,w\in W$，记
\[
 (\delta_xw)(y)=\begin{cases}w,&y=x,\\0,&y\ne x.\end{cases}
\]
\begin{definition}[项目中的紧性]
$P$ 紧，是指对每个 $v\in J$，存在有限集 $F_v\subseteq I$ 及映射 $U_{x,v}:W\to W$，使
\begin{equation}\label{eq:compact}
 P(f,v)=\sum_{x\in F_v}U_{x,v}(f(x))\qquad(f\in W^I).
\end{equation}
有限集与局部映射可以依赖 $v$；不要求一个统一的有限集合适用于全部 $v$。
\end{definition}

\begin{theorem}[紧性、有限依赖与无交互项]\label{thm:compact}
固定 $v$，记 $T(f)=P(f,v)$、$c=T(0)$。以下条件等价：
\begin{enumerate}
\item $T$ 具有\eqref{eq:compact}的形式。
\item 存在有限集 $F\subseteq I$，$T$ 只依赖 $f|_F$，且对支集不交的 $f,g$ 有
\begin{equation}\label{eq:orthadd}
 T(f+g)=T(f)+T(g)-c.
\end{equation}
\item $W$ 取离散拓扑、$W^I$ 取乘积拓扑时，$T$ 在 $0$ 处连续，且满足\eqref{eq:orthadd}。
\item 存在唯一的有限集 $E$ 和唯一的一族非零映射 $A_x:W\to W$，满足 $A_x(0)=0$，且
\begin{equation}\label{eq:canonical}
 T(f)=c+\sum_{x\in E}A_x(f(x)).
\end{equation}
\end{enumerate}
这里允许 $E=\varnothing$。规范局部映射为
\begin{equation}\label{eq:Ax}
 A_x(w)=T(\delta_xw)-c,
\end{equation}
并将恒零的项删去。对这类 $T$，在 $0$ 连续也等价于处处连续。
\end{theorem}
\begin{proof}
(1)$\Rightarrow$(2)：有限依赖显然。若 $\supp f\cap\supp g=\varnothing$，则每个坐标至少有一个取值为零，故
\[
 U_x(f(x)+g(x))=U_x(f(x))+U_x(g(x))-U_x(0).
\]
求和并用 $c=\sum_xU_x(0)$ 得\eqref{eq:orthadd}。

(2)$\Rightarrow$(4)：记 $f_F$ 为 $f$ 在 $F$ 上的截断。有限依赖给出 $T(f)=T(f_F)$。
对有限和 $f_F=\sum_{x\in F}\delta_xf(x)$ 逐次用\eqref{eq:orthadd}，得到
\[
 T(f)=c+\sum_{x\in F}\bigl(T(\delta_xf(x))-c\bigr).
\]
按\eqref{eq:Ax}定义 $A_x$，删去恒零项即得规范式。任一满足 $A_x(0)=0$ 的分解，在 $\delta_xw$ 处代入后都必须满足\eqref{eq:Ax}，故唯一。

(4)$\Rightarrow$(1)：若 $E\ne\varnothing$，把常数 $c$ 加入其中一项；若 $E=\varnothing$，任选 $x_0\in I$，令 $U_{x_0}\equiv c$ 即可。

(2)$\Rightarrow$(3)：固定有限个坐标的取值后输出不变，故 $T$ 处处连续。
反之，(3)意味着存在有限 $F$，只要 $g|_F=0$ 就有 $T(g)=c$。
将任意 $f$ 分解成支集不交的 $f_F+f_{I\setminus F}$，得到
\[
 T(f)=T(f_F)+T(f_{I\setminus F})-c=T(f_F).
\]
这就是有限依赖，再结合假设中的\eqref{eq:orthadd}得到(2)。
\end{proof}

\begin{remark}[不能删去无交互项条件]
有限依赖和连续性本身均不足以推出上述紧性。
例如 $W=\R$，取不同的 $a,b\in I$，令 $T(f)=f(a)f(b)$。
它只依赖两个坐标且连续，但
$T(\delta_a1+\delta_b1)=1$，而两项单独输入的输出都为零，因此不满足\eqref{eq:orthadd}。
\end{remark}

\begin{corollary}[加性差与线性差的紧性]\label{cor:lincompact}
若 $T:W^I\to W$ 可加，则以下条件等价：项目紧性；有限坐标依赖；离散乘积拓扑下在零点连续。
此时规范式中 $c=0$，且每个 $A_x$ 可加。若 $T$ 还是 $K$-线性的，则每个 $A_x\in\End_K(W)$。
\end{corollary}
\begin{proof}
可加性自动给出\eqref{eq:orthadd}，且 $T(0)=0$。
由 $A_x(w)=T(\delta_xw)$ 及 $\delta_x$ 的线性性立即得到相应的线性结论。
\end{proof}

\begin{corollary}[局部线性性无须另作假设]\label{cor:autolocal}
按修订后的定义，任何紧线性差都具有唯一的归一化表示
\[
 P(f,v)=\sum_{x\in F_v}T_{v,x}f(x),\qquad
 T_{v,x}\in\End_K(W),\quad T_{v,x}\ne0,
\]
其中 $F_v$ 有限，且 $T_{v,x}(w)=P(\delta_xw,v)$。
因此研究紧线性差时，无须再附加“存在一组线性的局部映射”这一假设。
\end{corollary}
\begin{proof}
规范分解的常数项为 $P(0,v)=0$；对 $w$ 的加性和齐次性均从 $P$ 继承。
唯一性来自定理\ref{thm:compact}。原始非归一化分解中的局部项可能带互相抵消的常数，故这里的结论针对归一化表示。
\end{proof}

\section{一般标量线性差：紧性之外的完整判据}\label{sec:noncompact}
本节只假设 $W=K$ 且 $P$ 为线性差，\textbf{暂不假设紧性}。
记 $E=K^I,F=K^J$，仍以 $\T:E\to F$ 表示 $P$，并令
\[
 \T^*:F^*\to E^*,\qquad \T^*(\Lambda)=\Lambda\circ\T,
\]
其中星号表示\textbf{完整代数对偶}。

\begin{theorem}[一般线性差的对偶存在唯一性]\label{thm:fullalgebraic}
方程 $\T f=h$ 可解，当且仅当
\begin{equation}\label{eq:fullcompat}
 \Lambda\in F^*,\quad \Lambda\circ\T=0
 \quad\Longrightarrow\quad \Lambda(h)=0.
\end{equation}
非空解集仍为 $f_0+\ker\T$。此外，每个右端的解都至多一个当且仅当 $\T^*$ 满射；每个右端可解当且仅当 $\T^*$ 单射。
\end{theorem}
\begin{proof}
若 $h=\T f$，必要性显然。若 $h\notin\im\T$，则商向量空间 $F/\im\T$ 中 $h$ 的像非零。
在含该向量的一组基上定义一个取值非零的泛函，再拉回到 $F$，就得到消灭 $\im\T$ 却不消灭 $h$ 的 $\Lambda$，违反\eqref{eq:fullcompat}。

若 $\T$ 单射，对任意 $\ell\in E^*$，先在 $\im\T$ 上定义 $\Lambda(\T f)=\ell(f)$，再延拓到 $F$，故 $\T^*$ 满射。
反之，若 $0\ne f\in\ker\T$，取 $\ell(f)\ne0$，则 $\ell$ 不可能属于 $\im\T^*$。
最后，$\ker\T^*$ 正是消灭 $\im\T$ 的泛函，故其为零当且仅当 $\im\T=F$。
\end{proof}

\begin{example}[有限关系相容仍不足：标量值反例]\label{ex:noncompact}
取 $K=V=\Q,n=1$，枚举 $\Q=\{q_1,q_2,\ldots\}$。
函数 $\delta_{q_j}1$ 与常值函数 $\mathbf1$ 构成线性无关族，将其延拓为 $\Q^{\Q}$ 的 Hamel 基。
令 $\ell_j$ 为对应于基向量 $\delta_{q_j}1$ 的坐标泛函，并定义
\[
 P(f,q_j)=\ell_j(f),\qquad h(q_j)=1.
\]
任意有限子系统可由相应的有限和 $f=\sum_{j\in F}\delta_{q_j}1$ 解出，且泛函族 $(\ell_j)$ 线性无关。
但任意函数作为 Hamel 基的线性组合只用有限项，故序列 $(\ell_j(f))_{j\ge1}$ 有限支集，不可能处处等于 $1$。
所以全局方程不可解。

每个 $\ell_j$ 都非有限取值型：若它是有限个点值的线性组合，对所有 $\delta_{q_k}1$ 代入会迫使它等于 $f\mapsto f(q_j)$；然而 $\ell_j(\mathbf1)=0$ 而 $\mathbf1(q_j)=1$，矛盾。
\end{example}
这里完整对偶判据没有失效。事实上 $\im\T=\Q^{(\Q)}$，在商 $\Q^{\Q}/\Q^{(\Q)}$ 中检测常值函数的泛函提供了障碍；该泛函不可能是有限个坐标评价的线性组合。
下一节的紧性恰好使所有存在性障碍都能由有限关系检测。

\section{标量值紧线性差：关系空间控制全部解}\label{sec:scalar}
设 $W=K$，$P$ 为紧的 $K$-线性差。由上一节，存在唯一系数 $a(v,x)$，每个固定 $v$ 只有有限个 $a(v,x)$ 非零，且
\begin{equation}\label{eq:rowfinite}
 P(f,v)=\sum_{x\in I}a(v,x)f(x).
\end{equation}
令
\[
 D_I=K^{(I)}=\bigoplus_{x\in I}Ke_x,\qquad
 D_J=K^{(J)}=\bigoplus_{v\in J}Ke_v.
\]
定义行向量、关系映射与行空间
\begin{equation}\label{eq:A}
 r_v=\sum_xa(v,x)e_x,\qquad
 A:D_J\to D_I,\ A(e_v)=r_v,\qquad L=\im A.
\end{equation}
每个函数 $f:I\to K$ 唯一对应线性泛函
$\ell_f:D_I\to K$，其中 $\ell_f(e_x)=f(x)$。
因此 $K^I=D_I^*$，且 $\T=A^*$。

\begin{theorem}[存在性的四个等价条件]\label{thm:scalar}
对 $h:J\to K$，以下等价：
\begin{enumerate}
\item 存在 $f:I\to K$ 使 $P(f,v)=h(v)$ 对全部 $v$ 成立。
\item 每个有限子方程组都可解。
\item 对任意有限族 $v_1,\ldots,v_s$ 与 $c_1,\ldots,c_s\in K$，都有
\begin{equation}\label{eq:compat}
 \sum_{i=1}^s c_ir_{v_i}=0
 \quad\Longrightarrow\quad
 \sum_{i=1}^s c_ih(v_i)=0.
\end{equation}
\item $\ell_h|_{\ker A}=0$，其中 $\ell_h(e_v)=h(v)$。
\end{enumerate}
\end{theorem}
\begin{proof}
(1)$\Rightarrow$(2)显然。若有限关系 $\sum_i c_ir_{v_i}=0$ 成立，取相关有限方程的一组解，将等式乘以 $c_i$ 相加，便得到\eqref{eq:compat}，故(2)$\Rightarrow$(3)。
由于 $D_J$ 的每个元素都是有限和，(3)与(4)完全相同。

由(3)，可在 $L$ 上定义
\[
 \varphi_h\!\left(\sum_i c_ir_{v_i}\right)=\sum_i c_ih(v_i).
\]
两个表达式的差是一条有限关系，所以定义良好且线性。
取 $L$ 的一组基并延拓为 $D_I$ 的基，在新增基向量上任意指定值，便将 $\varphi_h$ 延拓为 $\ell:D_I\to K$。
令 $f(x)=\ell(e_x)$，则 $P(f,v)=\ell(r_v)=h(v)$，得到(1)。
\end{proof}

\begin{theorem}[解空间、唯一性及任意右端可解]\label{thm:scalarstructure}
记 $C=D_I/L$。若 $f_0$ 为一组解，则
\begin{equation}\label{eq:scalarstructure}
 \Sol(h)=f_0+C^*,\qquad \ker\T\cong C^*.
\end{equation}
具体地，$C^*$ 通过商映射拉回到 $D_I^*$。此外：
\begin{enumerate}
\item 可解右端的解唯一，当且仅当 $L=D_I$，等价于每个 $e_x$ 都是有限个行向量的线性组合。
\item 每个 $h$ 都可解，当且仅当 $A$ 单射，等价于行族 $(r_v)_{v\in J}$ 线性无关。
\item 每个 $h$ 恰有一解，当且仅当行族 $(r_v)$ 是 $D_I$ 的一组基。
\end{enumerate}
\end{theorem}
\begin{proof}
齐次解对应恰好消灭 $L$ 的泛函，等价于 $C$ 上的任意泛函。
非零向量空间总存在非零线性泛函：把一个非零向量纳入基并在该向量上指定值 $1$。
所以 $C^*=0$ 当且仅当 $C=0$，得到(1)。
由定理\ref{thm:scalar}，所有 $h$ 可解等价于 $D_J$ 上所有泛函都消灭 $\ker A$，亦即 $\ker A=0$，得到(2)。结合两项得到(3)。
\end{proof}

\begin{corollary}[有限矛盾证书与有限恢复公式]
不可解时，存在一条有限关系\eqref{eq:compat}，其左侧为零而右侧非零。
若解唯一，对每个 $x$ 均可选有限表达式 $e_x=\sum_i c_ir_{v_i}$，于是
\[
 f(x)=\sum_i c_i h(v_i).
\]
不同表达式给出的值因相容性而相同。
\end{corollary}
\begin{proof}
若无解，定理\ref{thm:scalar}的相容条件必失败，故存在所述有限关系。
若解唯一，定理\ref{thm:scalarstructure}给出每个 $e_x$ 的有限行表示；
将解对应的泛函作用于该表示，即得恢复公式。
两个表示相减是一条有限行关系，所以相容性保证恢复值与表示无关。
\end{proof}

\begin{remark}[无限维数不能作形式减法]
若 $C$ 的一组基为 $B$，那么 $C^*\cong K^B$，全部自由参数是一族任意标量。
当 $|B|=d<\infty$ 时，齐次解空间维数为 $d$；当 $B$ 无限时，不可把 $K^B$ 当成直和 $K^{(B)}$，也不可写成未经解释的“未知数个数减秩”。
有限参数结构由商空间 $C$，而非由底层 $V$ 的维数直接决定。
\end{remark}

\begin{proposition}[精确序列与像的闭性]
存在精确序列
\[
 0\longrightarrow C^*\longrightarrow K^I
 \xrightarrow{\T}K^J\xrightarrow{\rho}(\ker A)^*
 \longrightarrow0,
\]
其中 $\rho(h)=\ell_h|_{\ker A}$。
赋予 $K$ 离散拓扑后，$\im\T$ 是 $K^J$ 中的闭子空间。
\end{proposition}
\begin{proof}
各处精确性来自前述延拓论证；$\rho$ 满射是因为 $\ker A$ 上的泛函可延拓到 $D_J$。
像由所有有限关系对应的方程 $\sum_i c_ih(v_i)=0$ 的交定义，每个此类方程的零集闭，故像闭。
\end{proof}

\section{有限维向量值情形的对偶判据}\label{sec:vectorfinite}
设 $0<\dim_KW=m<\infty$，并写
\[
 P(f,v)=\sum_x T_{v,x}f(x),\qquad T_{v,x}\in\End_K(W),
\]
每行仍只有有限个非零项。令
\[
 E_I=\bigoplus_{x\in I}W^*,\qquad E_J=\bigoplus_{v\in J}W^*.
\]
用 $i_x(\lambda)$ 表示位于 $x$ 坐标的 $\lambda\in W^*$，定义
\begin{equation}\label{eq:Bdual}
 B:E_J\to E_I,\qquad
 B(i_v(\lambda))=\sum_x i_x(\lambda\circ T_{v,x}).
\end{equation}
有限维性给出自然同构 $E_I^*\cong W^I$、$E_J^*\cong W^J$，且 $\T=B^*$。

\begin{theorem}[向量值存在唯一性]\label{thm:vectorfinite}
上述方程可解，当且仅当对任意有限族 $v_i\in J,\lambda_i\in W^*$，有
\begin{equation}\label{eq:veccompat}
 \sum_i\lambda_i\circ T_{v_i,x}=0\quad(\forall x\in I)
 \quad\Longrightarrow\quad
 \sum_i\lambda_i(h(v_i))=0.
\end{equation}
它也等价于每个有限子方程组可解。若可解，则
\[
 \Sol(h)=f_0+(E_I/\im B)^*.
\]
解唯一当且仅当 $B$ 满射；每个右端可解当且仅当 $B$ 单射。
\end{theorem}
\begin{proof}
在配对 $\langle f,i_x(\lambda)\rangle=\lambda(f(x))$ 下，直接计算得 $\T=B^*$。
条件\eqref{eq:veccompat}正是对应于 $h$ 的泛函消灭 $\ker B$。
同标量情形一样，先在 $\im B$ 上定义泛函，再延拓到 $E_I$；这里必须用 $W^{**}=W$ 才能把所得泛函解释为 $W$-值函数。

有限子方程可解蕴含\eqref{eq:veccompat}，因为每条关系只涉及有限个方程。
全局可解则显然蕴含有限可解。商空间、单射与满射的论证均与定理\ref{thm:scalarstructure}相同。
\end{proof}

\begin{remark}[标量系数与有限维值域是两种不同的充分条件]
若所有 $T_{v,x}=a(v,x)\id_W$，则即便 $W$ 无限维，定理\ref{thm:scalar}的证明仍可直接使用：
把 $\varphi_h:L\to W$ 延拓为 $D_I\to W$ 即可。
此时相容性是 $\sum_i c_ir_{v_i}=0\Rightarrow\sum_i c_i h(v_i)=0$，齐次解为
$\Hom_K(D_I/L,W)$。下一节的困难来自一般算子系数与无限维值域同时出现。
\end{remark}

\section{无限维值域：延拓障碍及明确反例}\label{sec:infinite}
\begin{example}[有限可解而全局不可解]\label{ex:infinite}
取 $V=\Q$、$n=1$、$K=\Q$、$W=\bigoplus_{j\ge1}\Q u_j$。
令 $\pi_j:W\to W$ 为 $\pi_j(\sum_k a_ku_k)=a_ju_1$。
在互异的参数 $v_j=j$ 处置
\[
 P(f,v_j)=\pi_j(f(0)),\qquad h(v_j)=u_1,
\]
其他参数处置 $P(f,v)=h(v)=0$。
该 $P$ 为紧的 $K$-线性差，每行甚至只取一个函数值。
任意有限子方程只要求 $f(0)$ 的有限个坐标等于 $1$，因而可解。
全部方程却要求 $f(0)=\sum_{j\ge1}u_j$，这不是 $W$ 的元素，故不可解。
\end{example}
该例说明不能在定理\ref{thm:vectorfinite}中删去有限维性。
若一般地使用 $W^*$ 作上述对偶延拓，只能先得到 $W^{**}$-值函数；
还需要保证各取值属于自然嵌入 $W\hookrightarrow W^{**}$ 的像。

\begin{theorem}[任意维值域的模同态判据]\label{thm:module}
设 $W\ne0$，$\calA=\End_K(W)$，以复合作为乘法，并把 $W$ 视为自然左 $\calA$-模。
令
\[
 \calM=\bigoplus_{x\in I}\calA e_x,\qquad
 \calN=\bigoplus_{v\in J}\calA e_v,
\]
定义左模同态
\[
 B_{\calA}:\calN\to\calM,\quad B_{\calA}(e_v)=\sum_x T_{v,x}e_x,
 \qquad\psi_h(e_v)=h(v).
\]
则 $P(f)=h$ 可解当且仅当下列两项同时成立：
\begin{enumerate}
\item $\psi_h(\ker B_{\calA})=0$，即对有限个 $a_i\in\calA$，
\[
 \sum_i a_iT_{v_i,x}=0\ (\forall x)\quad\Longrightarrow\quad
 \sum_i a_i h(v_i)=0;
\]
\item 由第一项诱导的左模同态
\[
 \bar\psi_h:\im B_{\calA}\to W,\qquad
 \bar\psi_h(B_{\calA}u)=\psi_h(u)
\]
可延拓为 $\calM\to W$ 的左模同态。
\end{enumerate}
有解时，齐次解空间自然为
\[
 \Hom_{\calA}(\coker B_{\calA},W).
\]
\end{theorem}
\begin{proof}
每个 $f:I\to W$ 唯一确定左模同态
$\phi_f:\calM\to W$，其中 $\phi_f(e_x)=f(x)$。
因此
\[
 (\phi_f B_{\calA})(e_v)=\sum_xT_{v,x}f(x)=P(f,v).
\]
所求方程等价于 $\psi_h=\phi_fB_{\calA}$。
先经像分解再延拓，恰好就是两项条件。差解满足 $\phi B_{\calA}=0$，因而恰好经余核分解。
\end{proof}
这里使用\textbf{左}模，所以相容式中的乘积顺序是 $a_iT_{v_i,x}$，不能交换。
在例\ref{ex:infinite}中，任何有限算子关系都满足第一项：取相关有限子系统的一个解即可验证。
因此该例也明确展示第二项可能失败，不能无证明地删去延拓条件。

\subsection{双对偶中的形式解与真实取值障碍}
继续使用 $E_I=\bigoplus_xW^*$、$E_J=\bigoplus_vW^*$ 和式\eqref{eq:Bdual}的 $B$，
但本小节允许 $W$ 无限维。记自然嵌入为 $j:W\hookrightarrow W^{**}$。

\begin{theorem}[有限相容性的双对偶解释]\label{thm:bidual}
对任意维 $W$，以下三个条件等价：
\begin{enumerate}
\item 每个有限子方程组 $\sum_xT_{v,x}f(x)=h(v)$ 都有 $W$-值解；
\item 式\eqref{eq:veccompat}的全部有限对偶关系成立；
\item 存在 $\widehat f:I\to W^{**}$，使
\begin{equation}\label{eq:bidualsolution}
 \sum_xT_{v,x}^{**}\widehat f(x)=j(h(v))\qquad(v\in J).
\end{equation}
\end{enumerate}
原来的 $W$-值问题有解，当且仅当可在第三项中选到 $\widehat f(x)\in j(W)$ 对所有 $x$ 成立。
\end{theorem}
\begin{proof}
第一项蕴含第二项，因为每条关系只涉及有限个方程。
反过来，固定有限参数集 $F\subset J$，令 $S\subset I$ 为这些行的支集并。
有限子系统是线性映射
$T_F:W^S\to W^F$ 的一个纤维。
若 $h|_F\notin\im T_F$，线性泛函分离给出 $\Lambda\in(W^F)^*$，消灭该像却不消灭 $h|_F$。
因为 $F$ 有限，$\Lambda$ 唯一可写成 $\sum_{v\in F}\lambda_v(w_v)$，其中 $\lambda_v\in W^*$。
消灭像恰表示 $\sum_{v\in F}\lambda_v\circ T_{v,x}=0$ 对全部 $x$ 成立，这与第二项矛盾。

第二项表示由 $h$ 给出的泛函 $\ell_h(i_v\lambda)=\lambda(h(v))$ 消灭 $\ker B$。
因此它先下降为 $\im B$ 上的泛函，再延拓为 $\ell\in E_I^*$。
自然识别 $E_I^*=(W^{**})^I$，令 $\widehat f(x)(\lambda)=\ell(i_x\lambda)$，
则 $\ell B=\ell_h$ 正好是\eqref{eq:bidualsolution}。
反过来，任何这样的形式解作用到有限关系上，都给出第二项。

若各 $\widehat f(x)$ 都属于 $j(W)$，由 $j$ 的单射性和 $T^{**}j=jT$，可唯一写成 $j(f(x))$ 并得到原方程。
原方程的解经 $j$ 嵌入又显然给出这种形式解。
\end{proof}

这说明有限维性的具体作用是 $j(W)=W^{**}$。
在无限维反例\ref{ex:infinite}中，形式解允许在全部坐标上取值 $1$，但该取值不能由 $W$ 的有限支集向量实现。
模同态延拓表述与这里的双对偶表述描述的是同一个全局取值障碍。

\section{模型方程：非齐次 Cauchy、Jensen 与差分多项式}\label{sec:models}
\subsection{非齐次 Cauchy 方程的完整相容条件}
设 $V,W$ 为 $\Q$-向量空间，维数不必有限。记
\[
 Df(x,y)=f(x+y)-f(x)-f(y).
\]
\begin{theorem}[对称二余循环判据]\label{thm:cocycle}
方程 $Df=h$ 有解，当且仅当
\begin{align}
 h(x,y)&=h(y,x),\label{eq:sym}\\
 h(x,y)+h(x+y,z)&=h(y,z)+h(x,y+z)\label{eq:cocycle}
\end{align}
对全部变量成立。有解时，全部解为
\[
 f_0+\Hom_{\Q}(V,W).
\]
若 $\dim_{\Q}V=d<\infty$、$\dim_{\Q}W=m<\infty$，则解仿射空间的维数为 $dm$。
\end{theorem}
\begin{proof}
必要性由展开 $Df$ 得到。证明充分性。
在\eqref{eq:cocycle}中令 $z=0$，得 $h(x+y,0)=h(y,0)$，故 $h(x,0)$ 为常值 $c=h(0,0)$。
置 $h_0=h-c$，则 $h_0(x,0)=h_0(0,x)=0$，且仍满足两式。

在集合 $E=W\times V$ 上定义加法
\[
 (u,x)+(w,y)=(u+w+h_0(x,y),x+y).
\]
对称性给出交换律，\eqref{eq:cocycle}给出结合律；$(0,0)$ 为零元，
$(-u-h_0(x,-x),-x)$ 为 $(u,x)$ 的逆元，故 $E$ 是阿贝尔群。
对每个正整数 $k$，设
\[
 b_k(x)=\sum_{j=1}^{k-1}h_0(jx,x),
\]
则归纳得 $k(u,x)=(ku+b_k(x),kx)$。
给定 $(w,y)$，只能取 $x=y/k$，再唯一取 $u=(w-b_k(x))/k$，便有 $k(u,x)=(w,y)$。
所以 $E$ 唯一可除，因而具有唯一的 $\Q$-向量空间结构。

投影 $E\to V$ 是满射群同态，因唯一可除性也是 $\Q$-线性映射。
选取一组基的提升，得到线性截面 $s:V\to E$，写为 $s(x)=(g(x),x)$。
等式 $s(x+y)=s(x)+s(y)$ 给出
$g(x+y)=g(x)+g(y)+h_0(x,y)$，即 $Dg=h_0$。
令 $f=g-c$，因为常值 $-c$ 的缺陷是 $c$，所以 $Df=h_0+c=h$。
最后，两个解之差可加，也即 $\Q$-线性；有限维结论由基的像可任意指定得到。
\end{proof}
上述证明同样适用于任意特征零域上的向量空间，但所得自由项一般仅为底层 $\Q$-线性映射。
不能在没有额外假设时将其换成 $\Hom_K(V,W)$。
指定一组 $\Q$-基上的函数值，可以从每个非空解仿射空间中选出唯一解。

\subsection{Jensen 方程与二次方程}
\begin{proposition}[有理向量空间上的两类齐次模型]
在 $V=\Q^d,W=\Q^m$ 上：
\begin{enumerate}
\item $2f((x+y)/2)-f(x)-f(y)=0$ 的全部解为 $f(x)=A(x)+b$，其中 $A$ 为 $\Q$-线性映射。
\item $f(x+y)+f(x-y)-2f(x)-2f(y)=0$ 的全部解为 $f(x)=B(x,x)$，其中 $B:V\times V\to W$ 为对称 $\Q$-双线性映射。
\end{enumerate}
解空间维数分别为 $m(d+1)$ 与 $m d(d+1)/2$。
\end{proposition}
\begin{proof}
第一项令 $g=f-f(0)$，先取 $y=0$ 得 $2g(x/2)=g(x)$，再代回得 $g(x+y)=g(x)+g(y)$。

第二项依次取 $x=y=0$、$x=0$、$y=x$，得到 $f(0)=0$、$f(-x)=f(x)$ 和 $f(2x)=4f(x)$。
定义
\[
 B(x,y)=\frac{f(x+y)-f(x-y)}4
       =\frac{f(x+y)-f(x)-f(y)}2.
\]
后一式表明对称性。分别对 $(x+y,z)$ 与 $(x-y,z)$ 应用原式并相减，得到
$B(x+z,y)+B(x-z,y)=2B(x,y)$。
由第一项且 $B(0,y)=0$，$B$ 对第一个变量线性；对称性给出双线性。
又 $B(x,x)=f(x)$。反向代入直接成立。维数按基上的独立参数计数。
\end{proof}

\subsection{高阶混合差分与普通多项式}
定义 $\Delta_hf(x)=f(x+h)-f(x)$。
\begin{theorem}[有理有限维情形的差分分类]\label{thm:polydifference}
设 $V=\Q^d,W=\Q^s$，$m\ge0$。以下等价：
\[
 \Delta_{h_1}\cdots\Delta_{h_{m+1}}f(x)=0
 \qquad(x,h_1,\ldots,h_{m+1}\in V);
\]
$f$ 的每个坐标是总次数不超过 $m$ 的普通有理系数多项式。
解空间维数为 $s\binom{d+m}{m}$。
\end{theorem}
\begin{proof}
一个差分使多项式总次数至少下降一次，故多项式必满足方程。
反过来，只需处理标量值。在整数格上，记 $e_i$ 为标准基、$\Delta_i=\Delta_{e_i}$。
对 $\alpha\in\N^d$ 置 $c_\alpha=(\Delta_1^{\alpha_1}\cdots\Delta_d^{\alpha_d}f)(0)$。
若 $|\alpha|>m$，混合差分假设给出 $c_\alpha=0$。
由移位算子 $T_i=1+\Delta_i$ 的二项式展开，对所有 $k\in\N^d$ 有
\[
 f(k)=\sum_{|\alpha|\le m}c_\alpha
       \prod_{i=1}^d\binom{k_i}{\alpha_i}=:p(k).
\]
右边是总次数不超过 $m$ 的多项式。沿各坐标轴，方程 $\Delta_i^{m+1}f=0$ 是首尾系数非零的线性递推，已知非负整数处的值后可向负方向唯一递推，故 $f=p$ 于 $\Z^d$。

对每个正整数 $N$，把同样的论证用于格 $N^{-1}\Z^d$，得到一个总次数不超过 $m$ 的多项式 $p_N$ 与 $f$ 在该格上一致。
在共同子格 $\Z^d$ 上有 $p_N=p$。多项式若在 $\Z^d$ 上为零，逐个变量使用一元多项式只有有限个根的事实，即知其恒零。
因此 $p_N=p$。每个有理点属于某个这样的格，所以 $f=p$ 于 $\Q^d$。
总次数不超过 $m$ 的单项式数为 $\binom{d+m}{m}$，得到维数公式。
\end{proof}
\begin{remark}
当定义域换成 $\R^d$ 而不加正则性时，普通实多项式分类不成立；
例如任意不连续加性函数都满足二阶差分为零。
一般应使用底层 $\Q$-结构上的广义多项式。本文只在已证明的有理有限维范围内使用上述普通多项式结论。
\end{remark}

\section{自然性之一：值域相容与群等变}\label{sec:natural}
“自然性”不是一个无需定义的单一条件。本节和下一节分别说明值域线性相容、固定群的等变性，以及跨全部空间的拉回自然性。

\subsection{与所有值域线性自映射交换}
\begin{proposition}[值域自然性强制标量系数]\label{prop:commutant}
设 $W\ne0$，$P$ 为紧的 $K$-线性差。若对所有 $B\in\End_K(W)$ 有
\[
 P(B\circ f,v)=B(P(f,v)),
\]
则每个局部算子 $T_{v,x}$ 都是标量倍数 $a(v,x)\id_W$。
逆命题亦成立，且不要求 $W$ 有限维。
\end{proposition}
\begin{proof}
取只在 $x$ 处非零的 $f$，得到 $T_{v,x}B=BT_{v,x}$。
以下证明任何与全部自同态交换的 $T$ 都是标量。
固定 $u\ne0$，取 $\lambda\in W^*$ 使 $\lambda(u)=1$，令 $p(w)=\lambda(w)u$。
因为 $Tp=pT$，所以 $Tu=pTu\in Ku$，写为 $Tu=au$。
对任意 $w$ 取线性映射 $B$ 使 $Bu=w$，则 $Tw=TBu=BTu=aw$，故 $T=a\id_W$。
反向显然。
\end{proof}
这是一条使无限维值域重新落入标量系数理论的自然条件，见第\ref{sec:vectorfinite}节末的说明。

\subsection{固定群作用下的行向量}
设群 $G$ 作用在集合 $I,J$ 上，并在 $D_I$ 上定义 $g e_x=e_{gx}$。
\begin{proposition}[等变性与行关系]
对标量值紧线性差，条件
\begin{equation}\label{eq:equiv}
 P(f\circ g^{-1},gv)=P(f,v)\qquad(g\in G)
\end{equation}
等价于
\[
 r_{gv}=g r_v.
\]
此时 $L=\Span\{r_v\}$ 是 $K[G]$-子模，$D_I/L$ 也自然成为 $K[G]$-模。
\end{proposition}
\begin{proof}
左边展开为 $\sum_xa(gv,x)f(g^{-1}x)=\sum_ya(gv,gy)f(y)$。
对所有 $f$ 相等，当且仅当 $a(gv,gy)=a(v,y)$，这正是行向量的等变式。
行空间的稳定性随即成立。
\end{proof}
因此，等变性把任意行空间变成一个带群作用的关系空间。但固定一个空间上的群等变性，仍弱于下一节对全部线性映射的自然性。

\subsection{非齐次右端的对称性与对称解}
采用函数上的作用 $g\cdot f=f\circ g^{-1}$，输出函数同样定义作用。
\begin{proposition}[解空间的仿射群作用]\label{prop:affinesolutions}
在式\eqref{eq:equiv}成立时，
\[
 g\cdot\Sol(h)=\Sol(g\cdot h).
\]
所以右端的稳定子 $H=\{g:g\cdot h=h\}$ 作用在非空解仿射空间上。
若解唯一，则该解必为 $H$-不变函数。
\end{proposition}
\begin{proof}
等变性等价于 $\T(g\cdot f)=g\cdot\T(f)$，因此 $\T(f)=h$ 蕴含 $\T(g\cdot f)=g\cdot h$；用 $g^{-1}$ 得到集合相等。
唯一解的每个 $H$-平移仍是解，故必等于自身。
\end{proof}

\begin{proposition}[对称解的精确障碍与有限群平均]\label{prop:symmetric}
若 $f_0\in\Sol(h)$，记 $Z=\ker\T$，并对 $g\in H$ 定义
$c(g)=g\cdot f_0-f_0$。
则 $c(g)\in Z$，且
\[
 c(gk)=g\cdot c(k)+c(g).
\]
存在 $H$-不变解，当且仅当存在 $u\in Z$ 使
\[
 c(g)=u-g\cdot u\qquad(g\in H).
\]
特别地，若 $H$ 有限且 $|H|$ 在 $K$ 中非零，任意可解的 $H$-不变右端都有 $H$-不变解。
\end{proposition}
\begin{proof}
$\T(c(g))=g\cdot h-h=0$，而将 $gk\cdot f_0-f_0$ 分成
$g\cdot(k\cdot f_0-f_0)+(g\cdot f_0-f_0)$ 得到所示恒等式。
全部解为 $f_0+u$；它不变等价于 $c(g)+g\cdot u-u=0$，即所述条件。
有限群情形可以直接取
$f_{\mathrm{av}}=|H|^{-1}\sum_{g\in H}g\cdot f_0$。
线性性给出 $\T(f_{\mathrm{av}})=h$，重新排列求和项给出不变性。
\end{proof}
上述 $c$ 是解仿射空间对应的一余循环；更换基点 $f_0$ 只使它增加一个余边界。
这里的恒等式及等价条件已经给出了所需的全部内容，不依赖上同调术语。

\begin{example}[特征整除群阶时平均结论会失败]
取 $K=V=\mathbb F_p$，平移群 $H=(\mathbb F_p,+)$，并令
$P(f,v)=\sum_{x\in\mathbb F_p}f(x)$、$h(v)=1$。
这是平移等变的紧线性差，$f=\delta_0 1$ 是解。
但平移不变函数只能为常数，而常数函数的输出为 $p c=0$，所以不存在平移不变解。
\end{example}

\subsection{平移等变性与群代数}
令 $G=(V,+)$，$J=V^n$ 采用对角平移。
将 $D_V$ 识别为群代数 $K[V]$，其基记作 $[x]$，乘法为 $[x][y]=[x+y]$。
\begin{theorem}[平移等变紧算子的群代数表述]\label{thm:translation}
若 $P$ 满足平移等变性，则它由一族有限群代数元素
\[
 b_d=r_{(0,d_2,\ldots,d_n)},\qquad d\in V^{n-1},
\]
唯一确定，且
\[
 r_{(v_1,\ldots,v_n)}=[v_1]b_{(v_2-v_1,\ldots,v_n-v_1)}.
\]
行空间恰是理想
\[
 L=(b_d:d\in V^{n-1})\subseteq K[V],
\]
齐次解自然为 $(K[V]/L)^*$。
特别地，$n=1$ 时恰为固定有限模板的卷积型算子
\[
 P(f,v)=\sum_{a\in F}c_a f(v+a).
\]
\end{theorem}
\begin{proof}
对每个元组减去首坐标，应用行等变式即得公式。
所有行正是所有 $[v]b_d$，其 $K$-线性包等于所示理想。
齐次解结论来自第\ref{sec:scalar}节。$n=1$ 时只剩一个基准行 $b$，其有限支集给出卷积模板。
\end{proof}

\begin{example}[周期与高阶差分]
周期方程 $f(x+a)-f(x)=0$ 对应理想 $([a]-[0])$。
商向量空间自然具有 $V/\Z a$ 的陪集基，因此解恰为各陪集上的任意常值。
更一般地，令增广理想
\[
 \mathfrak a=\ker\bigl(K[V]\to K,\ \sum c_x[x]\mapsto\sum c_x\bigr).
\]
$\mathfrak a$ 由所有 $[h]-[0]$ 生成，
$(m+1)$ 阶混合差分的行正是
\[
 [x]\prod_{i=1}^{m+1}([h_i]-[0]),
\]
其行理想为 $\mathfrak a^{m+1}$。故广义多项式空间为 $(K[V]/\mathfrak a^{m+1})^*$。
当 $K=\Q,V=\Q^d$ 时，第\ref{sec:models}节给出它的普通多项式解释和维数。
\end{example}

\section{固定空间上的线性群与自同态自然性}\label{sec:fixednatural}
本节设 $K$ 为无限域，$0<\dim_KV=d<\infty$，$W=K$。
先仅要求 $\mathrm{GL}(V)$-等变性，不要求在不同维数的空间间相容。
给定元组 $v=(v_1,\ldots,v_n)$，定义
\[
 a_v:K^n\to V,\quad a_v(\lambda)=\sum_i\lambda_iv_i,
 \qquad N_v=\ker a_v,\quad S_v=\im a_v.
\]
商映射诱导同构 $\bar a_v:K^n/N_v\to S_v$。

\begin{lemma}[有限支集必须位于输入张成空间]\label{lem:supportspan}
若 $P$ 紧且对 $\mathrm{GL}(V)$ 等变，则
\[
 \supp r_v\subseteq S_v.
\]
\end{lemma}
\begin{proof}
元组的稳定子由逐点固定 $S_v$ 的可逆线性变换组成；等变性使 $r_v$ 在此稳定子下不变。
若 $S_v=V$，结论显然。否则取补空间 $V=S_v\oplus U$。
任意 $x=s+u\notin S_v$ 满足 $u\ne0$。
令变换在 $S_v$ 上为恒等、在 $U$ 上乘以任意 $t\in K^\times$，得到无穷多个不同点 $s+tu$。
若 $x$ 处系数非零，整个无穷轨道上的系数均相同且非零，与行的有限支集矛盾。
\end{proof}

\begin{theorem}[固定空间的线性群等变分类]\label{thm:glclassification}
上述 $\mathrm{GL}(V)$-等变紧线性差，与下列数据一一对应：
对每个余维不超过 $d$ 的子空间 $N\le K^n$，任意选取有限支集元素
\[
 b_N\in K^{(K^n/N)}.
\]
对应的行向量是
\begin{equation}\label{eq:glclassification}
 r_v=(\bar a_v)_*b_{N_v},
\end{equation}
其中 $\bar a_v$ 的值看作 $V$ 的元素。不同 $N$ 的数据可独立选取。
\end{theorem}
\begin{proof}
两元组 $v,w$ 属于同一线性群轨道，当且仅当 $N_v=N_w$。
必要性显然；充分性由 $a_v(\lambda)\mapsto a_w(\lambda)$ 给出 $S_v\to S_w$ 的同构，再延拓为 $V$ 的自同构。
由引理\ref{lem:supportspan}，每个 $r_v$ 可沿 $\bar a_v$ 唯一拉回到 $K^{(K^n/N_v)}$。
同一轨道中的拉回相同，因为行向量满足 $r_{gv}=g_*r_v$，且稳定子逐点固定其支集。
所以得到唯一的 $b_N$。
反过来，任意这些数据通过\eqref{eq:glclassification}给出行有限算子；核空间在可逆变换下不变，公式直接保证等变性。
\end{proof}

\begin{remark}[为什么要求无限域]
在有限域和有限维 $V$ 上，所有轨道均有限。
例如 $P(f,v)=\sum_{x\in V}f(x)$ 对线性群等变，但在零元组处支集可以是整个 $V$，不包含于 $S_v=0$。
因此上述支集论证及其分类不能直接用于有限域；对全部空间及全部线性映射自然的分类则没有这一限制。
\end{remark}

\begin{theorem}[对全部线性自映射相容]\label{thm:endclassification}
在定理\ref{thm:glclassification}中进一步要求，对所有 $A\in\End_K(V)$ 有
\begin{equation}\label{eq:endnaturality}
 P(f\circ A,v)=P(f,Av).
\end{equation}
则这等价于：每当 $N\subseteq N'$ 且 $\operatorname{codim}N\le d$，都有
\begin{equation}\label{eq:quotientcompat}
 b_{N'}=(q_{N,N'})_*b_N,
 \qquad q_{N,N'}:K^n/N\to K^n/N'.
\end{equation}
特别地，若 $d\ge n$，全部数据由单个有限模板 $b_0\in K^{(K^n)}$ 唯一决定，因而
\[
 P(f,v)=\sum_{\lambda\in K^n}c_\lambda f\left(\sum_i\lambda_iv_i\right).
\]
\end{theorem}
\begin{proof}
式\eqref{eq:endnaturality}对全部 $f$ 成立，当且仅当 $r_{Av}=A_*r_v$。
总有 $N_v\subseteq N_{Av}$，且在两个商空间之间 $A$ 诱导的映射正是自然商映射，故\eqref{eq:quotientcompat}足以保证该行等式。

反过来，给定 $N\subseteq N'$，分别选取核为这两个空间的元组 $v,w$。
由包含关系，规则 $a_v(\lambda)\mapsto a_w(\lambda)$ 在 $S_v$ 上良定义且线性。
将其延拓到 $V$ 的自同态 $A$，便有 $Av=w$。
此时行等式恰好给出\eqref{eq:quotientcompat}。
若 $d\ge n$，可以取 $N=0$，所以 $b_N$ 都是 $b_0$ 的商推前；唯一性和公式立即成立。
\end{proof}

当 $d<n$ 时，条件\eqref{eq:quotientcompat}仍是完整的相容数据描述，但不能在此描述中凭空使用一个未出现的满秩元组来取 $b_0$。
下一节的跨空间自然性通过允许使用通用空间 $K^n$，消除了这一维数限制。

\section{自然性之二：对全部线性映射自然的分类}\label{sec:universal}
本节研究一整族算子 $P_V:K^V\times V^n\to K$，其中 $V$ 遍历全部 $K$-向量空间。
每个 $P_V$ 均为紧的 $K$-线性差。要求对任意线性映射 $A:V\to V'$、函数 $f:V'\to K$ 和元组 $v\in V^n$，有
\begin{equation}\label{eq:functorial}
 P_V(f\circ A,v)=P_{V'}(f,Av).
\end{equation}

\begin{theorem}[线性拉回自然性的有限模板分类]\label{thm:universal}
满足\eqref{eq:functorial}的算子族与有限支集系数族 $(c_\lambda)_{\lambda\in K^n}$ 一一对应，公式为
\begin{equation}\label{eq:universal}
 P_V(f,v_1,\ldots,v_n)
 =\sum_{\lambda\in K^n}c_\lambda
 f(\lambda_1v_1+\cdots+\lambda_nv_n).
\end{equation}
此处有限支集模板不依赖 $V$ 和 $v$。
\end{theorem}
\begin{proof}
在通用空间 $U=K^n$ 中取标准元组 $(e_1,\ldots,e_n)$。
紧线性性给出唯一有限行
$r=\sum_{\lambda\in K^n}c_\lambda[\lambda]\in K^{(K^n)}$。
对任意 $v$，存在线性映射 $A:U\to V$，$A(e_i)=v_i$。
在\eqref{eq:functorial}中使用此 $A$，得到
\[
 P_V(f,v)=P_U(f\circ A,e_1,\ldots,e_n)
 =\sum_\lambda c_\lambda f(A\lambda),
\]
这就是所示公式。通用元组处的行唯一，故系数唯一。
反过来，线性映射保持所有线性组合，故任意所示模板都满足自然性。
\end{proof}

\begin{corollary}[任意维值域的自然模板]\label{cor:operatornatural}
固定任意 $K$-向量空间 $W$。若紧线性差族
$P_V:W^V\times V^n\to W$ 对全部线性映射满足式\eqref{eq:functorial}，
则存在唯一有限支集算子族 $C_\lambda\in\End_K(W)$，使
\[
 P_V(f,v)=\sum_{\lambda\in K^n}C_\lambda
 \left(f\left(\sum_i\lambda_iv_i\right)\right).
\]
若此外与 $W$ 的全部线性自映射交换，且 $W\ne0$，则每个 $C_\lambda=c_\lambda\id_W$，回到标量系数模板。
\end{corollary}
\begin{proof}
在通用空间 $K^n$ 的标准元组处应用推论\ref{cor:autolocal}，得到有限族线性局部映射。
再沿 $e_i\mapsto v_i$ 拉回，逐字重复定理\ref{thm:universal}的证明即可。
值域相容时，在通用元组处逐坐标测试，得到 $C_\lambda$ 与全部自同态交换，再用命题\ref{prop:commutant}。
\end{proof}

\begin{corollary}[额外的平移与对称性条件]\label{cor:affine}
在定理\ref{thm:universal}的假设下：
\begin{enumerate}
\item 若还要求对所有空间的平移满足\eqref{eq:equiv}，则恰好只有
\[
 \lambda_1+\cdots+\lambda_n=1
\]
处允许 $c_\lambda\ne0$。也就是说，只能对仿射组合取值。
\item 对常值函数 $f\equiv b$ 有 $P_V(f,v)=b\sum_\lambda c_\lambda$。
因此消灭所有常数等价于 $\sum_\lambda c_\lambda=0$。
\item 对元组的任意置换不变，等价于模板系数在相应坐标置换下不变。
\end{enumerate}
\end{corollary}
\begin{proof}
第一项取 $V=K^{n+1}$、$v_i=e_i$、平移量 $t=e_{n+1}$。
平移前的行再整体平移，支点为 $(\lambda,1)$；在平移后的元组上应用公式，支点为 $(\lambda,\sum_i\lambda_i)$。
比较基向量系数，非零项必须满足坐标和为 $1$。充分性由仿射组合与平移交换直接得到。
第二项直接代入。第三项在通用元组处比较有限支集系数即可。
\end{proof}

\begin{example}[Cauchy 与 Jensen 的不同自然性]
Cauchy 差的模板为
$[(1,1)]-[(1,0)]-[(0,1)]$，故对全部线性映射自然。
它\textbf{不}满足对两个输入同时平移的等变性，也不消灭非零常值函数。
Jensen 差的模板为
$2[(1/2,1/2)]-[(1,0)]-[(0,1)]$（要求 $\operatorname{char}K\ne2$），
每个支点坐标和均为 $1$、系数总和为零，因此还具有平移等变性并消灭常数。
这说明不能把仿射自然性未经检查地加到所有 Cauchy 型关系上。
\end{example}

\begin{remark}[固定空间的线性群等变性较弱]
例如在任意非零 $V$ 上定义
$P(f,v)=b(v)f(0)$，其中 $b(0)=0$、$b(v)=1$ 对 $v\ne0$。
它是紧线性差，并对 $\mathrm{GL}(V)$ 等变。
但若将 $P$ 按此公式定义于所有空间，对零线性映射 $A:V\to V'$，
在 $v\ne0$ 处左边为 $f(0)$、右边为 $0$，故不满足\eqref{eq:functorial}。
因此定理\ref{thm:universal}的“全部线性映射”假设不能换成仅有线性自同构。
\end{remark}

\section{自然模板的矩与多项式解空间}\label{sec:moments}
本节设 $\operatorname{char}K=0$，$W=K$，并使用有限模板
\[
 P_V(f,v)=\sum_{\lambda\in F}c_\lambda f\left(\sum_i\lambda_iv_i\right),
 \qquad F\subset K^n\text{ 有限}.
\]
对多重指标 $\alpha=(\alpha_1,\ldots,\alpha_n)\in\N^n$，定义模板的矩
\[
 \mu_\alpha=\sum_{\lambda\in F}c_\lambda\lambda_1^{\alpha_1}\cdots\lambda_n^{\alpha_n}.
\]
其中零次幂按 $t^0=1$ 理解，包括 $t=0$。令
\begin{equation}\label{eq:symbol}
 \sigma_m(t_1,\ldots,t_n)=\sum_{\lambda\in F}c_\lambda
 (\lambda_1t_1+\cdots+\lambda_nt_n)^m
 =\sum_{|\alpha|=m}\binom m\alpha\mu_\alpha t^\alpha.
\end{equation}

\begin{theorem}[每个齐次次数上的消灭或单射二分]\label{thm:moment}
设 $0<\dim_KV=d<\infty$，记 $H_m(V)$ 为 $V$ 上 $m$ 次齐次多项式空间。
则下列条件等价：
\begin{enumerate}
\item $P_V$ 消灭 $H_m(V)$ 中的全部多项式；
\item 对全部 $|\alpha|=m$ 有 $\mu_\alpha=0$；
\item $\sigma_m=0$。
\end{enumerate}
若这三项不成立，则 $P_V$ 在 $H_m(V)$ 上是单射。
因此对任意 $M\ge0$，
\begin{equation}\label{eq:polykernel}
 \ker P_V\cap K[V]_{\le M}
 =\bigoplus_{\substack{0\le m\le M\\\sigma_m=0}}H_m(V),
\end{equation}
其维数为 $\sum_{0\le m\le M,\ \sigma_m=0}\binom{d+m-1}{m}$。
\end{theorem}
\begin{proof}
特征零时多项式展开式\eqref{eq:symbol}的多项式系数 $\binom m\alpha$ 非零，故第二、三项等价。
当 $m\ge1$ 时，把 $p\in H_m(V)$ 写成对称 $m$-线性形式的对角值 $p(x)=B(x,\ldots,x)$；具体可取
\[
 B(x_1,\ldots,x_m)=\frac1{m!}\sum_{S\subseteq\{1,\ldots,m\}}
 (-1)^{m-|S|}p\left(\sum_{i\in S}x_i\right).
\]
对每个齐次单项式展开时，容斥只保留每个 $x_i$ 都出现一次的项，因而该式对称且多线性，并在对角线上恢复 $p$。
多线性展开得
\[
 P_V(p,v)=\sum_{|\alpha|=m}\binom m\alpha\mu_\alpha
 B(\underbrace{v_1,\ldots,v_1}_{\alpha_1},\ldots,
   \underbrace{v_n,\ldots,v_n}_{\alpha_n}).
\]
所以第二项蕴含第一项。$m=0$ 时直接用常数函数验证同一结论。

若 $\sigma_m\ne0$，由于 $K$ 无限，可取 $t\in K^n$ 使 $\sigma_m(t)\ne0$。
对每个 $x\in V$，齐次性给出
\begin{equation}\label{eq:homorecovery}
 P_V(p,t_1x,\ldots,t_nx)=\sigma_m(t)p(x).
\end{equation}
因此 $P_V(p)=0$ 强制 $p=0$，说明单射性，也说明不能消灭非零的 $H_m(V)$。

最后把 $p=\sum_{m=0}^M p_m$ 按次数分解。
若 $P_V(p)=0$，将输入元组同时乘以任意 $s\in K$，得到
$\sum_m s^mP_V(p_m,v)=0$。
无限域上一元多项式系数唯一，故每个 $P_V(p_m,v)=0$。
应用前述二分即得\eqref{eq:polykernel}，维数由齐次单项式计数给出。
\end{proof}

\begin{corollary}[非齐次多项式解的恢复与检验]\label{cor:polyrecover}
若要求 $f\in K[V]_{\le M}$，则可解右端 $h$ 必须是在元组变量中总次数不超过 $M$ 的多项式。
写其齐次分解为 $h=\sum_{m=0}^M h_m$。
对每个 $\sigma_m=0$，必须有 $h_m=0$，此时 $f_m\in H_m(V)$ 任意。
对每个 $\sigma_m\ne0$，选取 $t$ 使 $\sigma_m(t)\ne0$，唯一候选为
\[
 f_m(x)=\frac{h_m(t_1x,\ldots,t_nx)}{\sigma_m(t)}.
\]
方程可解当且仅当这些候选逐个满足 $P_V(f_m)=h_m$。
\end{corollary}
\begin{proof}
必要性与候选公式来自次数保持和\eqref{eq:homorecovery}。
若各次候选均通过代入检验，将它们与任意被消灭次数的齐次项相加便得到解。
\end{proof}
这里的最终代入检验不能省去：仅在全部变量共线的元组上恢复候选，并不自动保证对一般元组的方程成立。

\begin{proposition}[有限模板的最低非零矩有界]\label{prop:momentbound}
若模板非零，合并重合支点并删去零系数后共有 $s$ 个支点，则存在 $|\alpha|\le s-1$ 使 $\mu_\alpha\ne0$。
因此非零有限模板不能消灭所有多项式。
\end{proposition}
\begin{proof}
固定一个系数非零的支点 $\lambda_0$。
对每个其余支点 $\lambda$，选取一个坐标 $i$ 使 $\lambda_i\ne(\lambda_0)_i$，并取仿射一次多项式
$L_\lambda(z)=(z_i-\lambda_i)/((\lambda_0)_i-\lambda_i)$。
乘积 $q=\prod_{\lambda\ne\lambda_0}L_\lambda$ 次数至多 $s-1$，在 $\lambda_0$ 处为 $1$，在其余支点处为 $0$。
故 $\sum_\lambda c_\lambda q(\lambda)=c_{\lambda_0}\ne0$。
把 $q$ 展开为单项式线性组合，便有至少一个相应的矩非零。
\end{proof}

\begin{example}[Cauchy、Jensen 与差分的次数结构]
Cauchy 差的符号为 $\sigma_m(t,u)=(t+u)^m-t^m-u^m$。
其中 $\sigma_0=-1$、$\sigma_1=0$，而 $m\ge2$ 时混合项非零，所以多项式核恰为一次齐次多项式。
Jensen 差的符号为 $2((t+u)/2)^m-t^m-u^m$；它恰在 $m=0,1$ 时为零，所以多项式核恰为仿射函数。

对 $r\ge1$，重复步长差分
\[
 \Delta_h^r f(x)=\sum_{k=0}^r(-1)^{r-k}\binom rk f(x+kh)
\]
有 $r+1$ 个模板支点 $(1,k)$。
其次数小于 $r$ 的矩全为零，而
$\mu_{(m-r,r)}=\sum_k(-1)^{r-k}\binom rk k^r=r!\ne0$ 对每个 $m\ge r$ 成立。
该恒等式来自一元 $r$ 阶差分将 $t^r$ 化为常数 $r!$。
故其多项式核恰为次数小于 $r$ 的多项式，且命题\ref{prop:momentbound}的界可达到。
\end{example}
上述结论只是在多项式函数类内部的分类。对任意函数求全局核时，不得据此排除不连续加性函数等非多项式解。

\section{一元自然算子：乘法群代数与完整解判据}\label{sec:dilation}
现在取 $n=1$、$V\ne0$、$W=K$，考虑自然模板
\begin{equation}\label{eq:dilation}
 P(f,v)=\sum_{\lambda\in K}c_\lambda f(\lambda v),
\end{equation}
系数有限支集。本节的群代数判据适用于任意域。
记
\[
 s=\sum_{\lambda\in K}c_\lambda,\qquad
 \mathcal R=K[K^\times],\qquad
 b=\sum_{\lambda\in K^\times}c_\lambda[\lambda]\in\mathcal R.
\]
每条一维子空间 $L\subseteq V$ 选非零代表 $u_L$。

\begin{theorem}[按射影直线分解非齐次问题]\label{thm:dilationsolve}
方程 $P(f)=h$ 有解，当且仅当可选择 $z\in K$，使 $sz=h(0)$，并且对每条 $L$，函数
\[
 k_L(t)=h(tu_L)-c_0z\qquad(t\in K^\times)
\]
对应的线性泛函 $\ell_{k_L}:\mathcal R\to K$ 消灭
\[
 \operatorname{Ann}_{\mathcal R}(b)=\{a\in\mathcal R:ba=0\}.
\]
固定这样的 $z=f(0)$ 后，各条 $L$ 上的解相互独立；每一条上的自由解空间为 $(\mathcal R/b\mathcal R)^*$。
\end{theorem}
\begin{proof}
原点方程为 $sf(0)=h(0)$。
对 $v=tu_L\ne0$，置 $g_L(t)=f(tu_L)$，方程变为
\[
 \sum_{\lambda\ne0}c_\lambda g_L(\lambda t)=k_L(t).
\]
将 $g_L$ 识别为 $\mathcal R$ 上的线性泛函，此算子就是乘法映射
$M_b:\mathcal R\to\mathcal R,\ a\mapsto ba$ 的代数对偶。
第\ref{sec:scalar}节的关系空间判据给出相容条件 $\ell_{k_L}|_{\ker M_b}=0$，自由解为余核的对偶。
所有非零向量恰好属于一条这样的 $L$，且式\eqref{eq:dilation}不混合不同 $L$，所以这些解可以独立拼接。
\end{proof}

\begin{corollary}[一元自然算子的整体存在唯一性]\label{cor:dilationbij}
对式\eqref{eq:dilation}有：
\begin{enumerate}
\item 算子 $\T:K^V\to K^V$ 单射，当且仅当 $s\ne0$ 且 $b$ 为 $\mathcal R$ 中的单位；在这种情形下它自动为双射。
\item 每个 $h$ 都可解，当且仅当 $s\ne0$ 且 $M_b$ 单射，即 $b$ 的零化子为零。
\item 因此，至多一解与对任意右端恰有一解，对这类一元自然算子等价。
\end{enumerate}
\end{corollary}
\begin{proof}
若 $s=0$，所有常值函数都在核中，且非零 $h(0)$ 无法实现，所以单射和满射都不成立。
设 $s\ne0$。齐次解的原点值必须为零；其余自由项为各条 $L$ 上的 $(\mathcal R/b\mathcal R)^*$。
因此核为零当且仅当 $b\mathcal R=\mathcal R$，这又等价于 $b$ 为单位；此时每条直线上的乘法映射可逆，所以所有右端均可唯一求解。
任意右端可解则等价于 $M_b^*$ 满射，按代数对偶判据等价于 $M_b$ 单射。
\end{proof}

\begin{example}[反射模板的显式分类]
设 $\operatorname{char}K\ne2$，考虑 $P(f,v)=f(v)-a f(-v)$。
若 $a^2\ne1$，则任意右端均有唯一解
\[
 f(v)=\frac{h(v)+a h(-v)}{1-a^2}.
\]
这由原方程与在 $-v$ 处的方程联立得到，代入也可直接验证。
若 $a=1$，可解当且仅当 $h(-v)=-h(v)$，全部解为 $h/2$ 加任意偶函数。
若 $a=-1$，可解当且仅当 $h(-v)=h(v)$，全部解为 $h/2$ 加任意奇函数。
\end{example}

\begin{example}[任意右端可解而不唯一]
在特征零域上取 $P(f,v)=f(2v)-2f(v)$。
原点处强制 $f(0)=-h(0)$。
在非零向量集合上按 $v\sim2^k v$ 分成轨道，每条轨道选一个代表 $u$，任意指定 $f(u)$，然后用
\[
 f(2^{k+1}u)=2f(2^ku)+h(2^ku)
\]
向正、负两个方向递推。
特征零保证轨道没有周期，因此得到全部解。任意 $h$ 均可解，每条非零轨道各有一个自由参数。
\end{example}



\section{一般自然模板的单射性与有限域模型}\label{sec:generalideal}
本节仍取标量值 $W=K$，并固定完全线性自然模板
\[
 P_V(f,v)=\sum_{\lambda\in K^n}c_\lambda
 f\!\left(\sum_i\lambda_i v_i\right),\qquad
 s=\sum_\lambda c_\lambda.
\]
模板有限支集，但域 $K$ 及空间 $V$ 不要求有限。
以下关于非零解和单射性的统一判据均假设 $V\ne0$。

\subsection{单射性只需检验一维空间}
\begin{theorem}[一维检测原则]\label{thm:onedetect}
对任意非零 $K$-空间 $V$，
\[
 \T_V\text{ 单射}\quad\Longleftrightarrow\quad\T_K\text{ 单射}.
\]
若 $s=0$，二者都不单射。
\end{theorem}
\begin{proof}
若 $\T_K$ 单射，设 $f\in\ker\T_V$。对每个 $u\in V$，线性映射
$A_u:K\to V,\ t\mapsto tu$ 的拉回自然性给出
$f\circ A_u\in\ker\T_K$，所以 $f(u)=0$，从而 $f=0$。

反过来，若 $0\ne g\in\ker\T_K$，选择一个满射线性泛函
$\ell:V\to K$。自然性给出 $g\circ\ell\in\ker\T_V$，而满射性保证它非零。
最后 $s=0$ 时任意常值函数都属于齐次核。
\end{proof}

令 $\mathscr R=K[K^\times]$。对 $t=(t_1,\ldots,t_n)\in K^n$，定义
\[
 b_t=\sum_{\lambda\cdot t\ne0}c_\lambda[\lambda\cdot t],
 \qquad
 I_c=(b_t:t\in K^n)\subseteq\mathscr R.
\]
其中相同的 $\lambda\cdot t$ 所对应的系数要合并。由
$b_{at}=[a]b_t$（$a\ne0$），这些 $b_t$ 的线性包已经是理想 $I_c$。

\begin{theorem}[乘法群代数的理想判据]\label{thm:generalideal}
在一维空间 $K$ 上有自然的向量空间同构
\[
 \ker\T_K\cong \{k\in K:sk=0\}\oplus(\mathscr R/I_c)^*.
\]
因此对任何 $V\ne0$，
\begin{equation}\label{eq:generalinject}
 \T_V\text{ 单射}
 \quad\Longleftrightarrow\quad s\ne0\ \text{且}\ I_c=\mathscr R.
\end{equation}
\end{theorem}
\begin{proof}
齐次解在零元组处满足 $sf(0)=0$。写 $f=k+g$，其中
$k=f(0)$、$g(0)=0$。因为 $sk=0$，有 $\T_Kf=0$ 当且仅当
$\T_Kg=0$。将 $g|_{K^\times}$ 识别为 $\mathscr R$ 上的线性泛函，
$\T_Kg(t)=0$ 正是该泛函消灭 $b_t$ 的条件，故等价于消灭 $I_c$。
这给出分解。非零向量空间有非零对偶，再应用定理\ref{thm:onedetect}，得到单射性判据。
\end{proof}

\begin{remark}
式\eqref{eq:generalinject}是所有非零维数上的单射性判据；
当 $n\ge2$ 时，它并未声称 $\ker\T_V$ 可以按射影直线独立指定。
不同方向的输入会出现在同一个方程中，通常产生额外耦合。
第\ref{sec:dilation}节的逐直线独立分解只适用于一元伸缩模板。
\end{remark}

\subsection{有限域上的字符判据}
设 $K=\mathbb F_q$。其乘法群循环，阶为 $q-1$，并且 $q-1$ 在 $K$ 中非零。
对 $1\le d\le q-1$，定义函数
\[
 E_d(t)=\sum_\lambda c_\lambda(\lambda\cdot t)^d
 \qquad(t\in K^n).
\]
这里 $0^d=0$，不使用 $0^0$。

\begin{proposition}[有限域的一维齐次核]\label{prop:finitechars}
有
\[
 \ker\T_K
 =
 \Span_K\!\left(
 \{1:s=0\}\ \cup\
 \{\,x\mapsto x^d:1\le d\le q-1,\ E_d\equiv0\,\}
 \right),
\]
且显示出的函数构成一组基。因此
\[
 \dim_K\ker\T_K=\mathbf1_{s=0}
 +\#\{d\in\{1,\ldots,q-1\}:E_d\equiv0\}.
\]
对任意 $V\ne0$，$\T_V$ 单射当且仅当 $s\ne0$ 且每个 $E_d$ 都不是零函数。
\end{proposition}
\begin{proof}
取 $K^\times$ 的生成元 $\omega$，则
\[
 \mathscr R\cong K[z]/(z^{q-1}-1).
\]
多项式 $z^{q-1}-1$ 的根恰为 $K^\times$，且全为单根。
因此通过在这些根处取值，有
$\mathscr R\cong\prod_{d=1}^{q-1}K$。
各坐标可写为群代数字符
$[\lambda]\mapsto\lambda^d$；$d=q-1$ 是平凡字符。
理想 $I_c$ 在第 $d$ 个坐标为零，当且仅当所有 $b_t$ 在该字符下取值为零，
即 $E_d(t)=0$ 对全部 $t$ 成立。
该坐标上的对偶函数恰为 $x^d$ 在 $K^\times$ 上的限制，并在 $0$ 处补零。
再加上前一定理中的常数项即得结论。
这些函数也就是不同的约化单项式 $1,x,\ldots,x^{q-1}$ 的一部分，故线性无关。
\end{proof}

\subsection{有限维有限域的全部函数解：约化单项式模型}
设 $V=K^d$，$d\ge1$。每个函数 $f:K^d\to K$ 唯一写成
\[
 f(x)=\sum_{\beta\in\{0,\ldots,q-1\}^d}a_\beta x^\beta.
\]
存在性可由逐坐标 Lagrange 插值得到；唯一性来自逐变量使用
“次数小于 $q$ 的非零多项式至多有其次数个根”。

为避免把正特征的二项式系数当成非零数，保留如下完整系数。
对非负整数矩阵 $\Gamma=(\gamma_{i\ell})\in\N^{n\times d}$，置
\[
 \beta_\ell(\Gamma)=\sum_{i=1}^n\gamma_{i\ell},\qquad
 \rho_i(\Gamma)=\sum_{\ell=1}^d\gamma_{i\ell},\qquad
 \mu_\rho=\sum_\lambda c_\lambda\prod_i\lambda_i^{\rho_i}.
\]
在 $\mu_0$ 中约定 $0^0=1$，从而 $\mu_0=s$。

\begin{theorem}[有限域全核的有限系数判据]\label{thm:finitefull}
单项式 $x^\beta$ 属于 $\ker\T_V$，当且仅当对每个满足
$\beta(\Gamma)=\beta$ 的矩阵 $\Gamma$，都有
\begin{equation}\label{eq:finitecoef}
 \left(\prod_{\ell=1}^d
 \binom{\beta_\ell}{\gamma_{1\ell},\ldots,\gamma_{n\ell}}\right)
 \mu_{\rho(\Gamma)}=0\quad\text{于 }K.
\end{equation}
全部通过该判据的约化单项式构成 $\ker\T_V$ 的一组基。
\end{theorem}
\begin{proof}
将输入元组写成 $v_i=(v_{i1},\ldots,v_{id})$。多项式展开给出
\[
 \T_V(x^\beta)(v)
 =
 \sum_{\beta(\Gamma)=\beta}
 \left(\prod_{\ell=1}^d
 \binom{\beta_\ell}{\gamma_{1\ell},\ldots,\gamma_{n\ell}}\right)
 \mu_{\rho(\Gamma)}
 \prod_{i,\ell}v_{i\ell}^{\gamma_{i\ell}}.
\]
每个 $\gamma_{i\ell}\le\beta_\ell\le q-1$，所以右边已是约化多项式。
它表示零函数当且仅当全部系数为零，得到判据。
若 $\beta\ne\beta'$，对应输出单项式在某一列的总次数不同，故没有重合项。
因此不同 $\beta$ 的非零输出之间不可能抵消；核正是显示出的单项式张成空间。
\end{proof}

非齐次方程可直接列成 $q^{nd}\times q^d$ 的取值矩阵系统；
也可用上述系数展开建立等价系统，逐个输出单项式比较系数。
其有限矩阵求解给出相容条件、一个特解和核的基。
有限域上的“多项式表示全部函数”依赖约化次数约定，不能移植到特征零域来排除任意函数解。

\section{特征零的秩二低支撑模板与方向消元}\label{sec:lowrank}
从本节到第\ref{sec:normexample}节，除特别说明外统一假设
\[
 \operatorname{char}K=0,\qquad W=K,\qquad V\ne0.
\]
自然模板的系数支集记作 $\Lambda\subset K^n$，所有重合支点已经合并，
零系数已经删去。

\subsection{模板秩的约化与非齐次右端}
若 $\Span_K\Lambda$ 的维数为 $d_0$，取其一组基
$p_1,\ldots,p_{d_0}\in K^n$，写
$\lambda=\sum_j a_{\lambda j}p_j$，并令
\[
 x_j=\sum_i(p_j)_i v_i.
\]
于是原方程通过满射线性映射
\[
 \pi:V^n\longrightarrow V^{d_0},\qquad v\longmapsto(x_1,\ldots,x_{d_0})
\]
分解。满射性来自系数矩阵满行秩，因而有一个 $K$-线性右逆。
齐次问题与约化后的问题等价。
对非齐次问题，必须且只需先要求 $h=\widetilde h\circ\pi$，
再求解约化方程 $P_{\mathrm{red}}f=\widetilde h$。
若 $h$ 在某条 $\pi$ 的纤维上不恒定，原方程立即不可解。

秩零模板只含可能的原点项，即 $P(f,v)=c_0f(0)$；
秩一经上述约化归入第\ref{sec:dilation}节。
以下固定秩为二并要求总支撑不超过四点。

写成
\begin{equation}\label{eq:ranktwo}
 P(f,x,y)=c_0f(0)+\sum_{j=1}^m c_jf(a_jx+b_jy),
 \qquad x,y\in V,
\end{equation}
其中 $c_j\ne0$，$(a_j,b_j)\ne(0,0)$ 两两不同并张成 $K^2$，且
\[
 m+\mathbf1_{c_0\ne0}\le4.
\]
记
\[
 s=c_0+\sum_{j=1}^m c_j,\qquad
 r=\#\{K(a_j,b_j):1\le j\le m\}.
\]
这里 $r$ 是过原点的方向数，必有 $2\le r\le4$。
原点支撑计入总支撑数，但不计入 $r$。

\begin{center}
\begin{tabular}{p{.40\linewidth}p{.50\linewidth}}
\toprule
非零支撑的方向分布 & 全部齐次解的可能形式\\
\midrule
两个方向，其中至少一个只有一点 & 常值函数 $k$，且 $sk=0$\\
两个方向，每个恰有两点 & 常数加由两个伸缩律控制的函数\\
三个方向 & 常数加 $\Q$-线性函数\\
四个方向 & 常数加 $\Q$-线性函数加对称 $\Q$-双线性型的对角值\\
\bottomrule
\end{tabular}
\end{center}

\subsection{单点方向的差分消元}
\begin{theorem}[方向数控制差分次数]\label{thm:singleton}
对式\eqref{eq:ranktwo}，不必先假设支撑至多四点。
若共有 $r\ge2$ 个非零方向，且其中一个方向只有一个支撑点，
则每个齐次解都满足
\begin{equation}\label{eq:directiondifference}
 \Delta_{h_1}\cdots\Delta_{h_{r-1}}f(z)=0
 \qquad(z,h_1,\ldots,h_{r-1}\in V).
\end{equation}
\end{theorem}
\begin{proof}
把同一方向上的项合并为一个函数，原方程可写成
\[
 c_0f(0)+\sum_{\nu=1}^rF_\nu(L_\nu(x,y))=0,
\]
其中 $L_\nu$ 为两两不成比例的非零线性形式。
对单点方向可写 $F_1(z)=cf(tz)$，其中 $c,t\in K^\times$。

对每个 $\nu=2,\ldots,r$，任给 $h_\nu\in V$。
$L_1,L_\nu$ 的系数矩阵可逆，故存在 $\eta_\nu\in V^2$ 使
$L_\nu(\eta_\nu)=0$、$L_1(\eta_\nu)=h_\nu$。
对参数 $(x,y)$ 依次施加平移差分
$\Delta_{\eta_2}\cdots\Delta_{\eta_r}$。
第 $\nu$ 个方向的项被 $\Delta_{\eta_\nu}$ 消去；
所有差分彼此交换，常数项也被消去，故最后仅剩
\[
 \Delta_{h_2}\cdots\Delta_{h_r}F_1(L_1(x,y))=0.
\]
$L_1$ 满射，$t$ 非零，且各 $h_\nu$ 任意；除以 $c$ 并重新命名步长即得结论。
\end{proof}

该证明解释了单点方向的作用：重复方向只能先合并成未知函数，
只有单点方向能在消元结束后直接恢复原来的 $f$。
相关复数域线性函数方程对孤立方向的研究可参见\cite{KissVarga}；
这里的差分证明独立成立，并不使用该文的其他分类结果。

\begin{lemma}[零至二次广义多项式的唯一分解]\label{lem:lowdegree}
对 $\Q$-向量空间 $V$ 和特征零值域 $K$：
\begin{enumerate}
\item 一阶差分恒零，当且仅当 $f$ 为常数；
\item 二阶混合差分恒零，当且仅当 $f=k+A$，$A:V\to K$ 为 $\Q$-线性映射；
\item 三阶混合差分恒零，当且仅当
\[
 f(x)=k+A(x)+B(x,x),
\]
其中 $A$ 为 $\Q$-线性映射，$B$ 为对称 $\Q$-双线性映射。
\end{enumerate}
上述分解均唯一。
\end{lemma}
\begin{proof}
第一项直接令基点为零。第二项取 $k=f(0)$，
在零点处写出 $\Delta_u\Delta_vf(0)=0$，就得到 $f-k$ 可加，
而可加映射在 $\Q$-向量空间间自动为 $\Q$-线性。

证明第三项。令 $g=f-f(0)$ 并定义
\[
 B(u,v)=\frac12\Delta_u\Delta_vg(0)
 =\frac12\bigl(g(u+v)-g(u)-g(v)\bigr).
\]
三阶差分为零使二阶差分与基点无关。
利用 $\Delta_{u+v}=\Delta_u+T_u\Delta_v$，其中 $T_u$ 是平移算子，得到
\[
 \Delta_{u+v}\Delta_wg(0)
 =\Delta_u\Delta_wg(0)+\Delta_v\Delta_wg(u)
 =\Delta_u\Delta_wg(0)+\Delta_v\Delta_wg(0).
\]
因此 $B$ 对第一变量可加；对称性给出第二变量可加，进而对两变量 $\Q$-线性。
令 $A(x)=g(x)-B(x,x)$。由
\[
 g(u+v)-g(u)-g(v)=2B(u,v)
\]
及双线性展开，$A$ 可加。反向直接差分验证。
唯一性可从
\[
 k=f(0),\qquad
 A(x)=\frac{f(x)-f(-x)}2,\qquad
 B(x,x)=\frac{f(x)+f(-x)}2-f(0)
\]
以及对称双线性型的极化公式得到。
\end{proof}

\begin{corollary}[四点范围内的次数上界]\label{cor:degreebound}
在总支撑不超过四点、秩二的假设下：
若 $r=2$ 且有单点方向，全部解恰为满足 $sk=0$ 的常数；
若 $r=3$，全部解至多为常数加加性项；
若 $r=4$，全部解至多再含一个二次广义多项式项。
\end{corollary}
\begin{proof}
$r=3$ 或 $4$ 时，由抽屉原理必有单点方向。
结合定理\ref{thm:singleton}和引理\ref{lem:lowdegree}即可。
$r=2$ 时差分结论直接迫使常数，代入给出 $sk=0$。
\end{proof}

\section{三个、四个方向的精确核与二方向例外}\label{sec:kernelclassification}
\subsection{加性项与二次项的充要条件}
定义
\[
 \begin{split}
 \mathcal A_c(V)=\{A\in\Hom_{\Q}(V,K):\
 &\sum_j c_jA(a_jx)=0,\\
 &\sum_j c_jA(b_jx)=0\quad(\forall x\in V)\}.
 \end{split}
\]
定义 $\mathcal B_c(V)$ 为所有对称 $\Q$-双线性型
$B:V\times V\to K$ 中满足
\begin{equation}\label{eq:quadraticconstraints}
 \begin{aligned}
 \sum_jc_jB(a_ju,a_jv)&=0,\\
 \sum_jc_jB(b_ju,b_jv)&=0,\\
 \sum_jc_jB(a_ju,b_jv)&=0
 \end{aligned}
 \qquad(u,v\in V)
\end{equation}
者。

\begin{theorem}[三个、四个方向的完整齐次解]\label{thm:exactlowkernel}
在式\eqref{eq:ranktwo}及四点支撑假设下：
\begin{align}
 r=3:\quad
 \ker\T_V&=\{\,k+A:sk=0,\ A\in\mathcal A_c(V)\,\};
 \label{eq:threekernel}\\
 r=4:\quad
 \ker\T_V&=\{\,x\mapsto k+A(x)+B(x,x):
 sk=0,\ A\in\mathcal A_c(V),\ B\in\mathcal B_c(V)\,\}.
 \label{eq:fourkernel}
\end{align}
各参数唯一。
\end{theorem}
\begin{proof}
次数上界由推论\ref{cor:degreebound}给出。
先将 $(x,y)$ 同时乘以任意 $t\in\Q$。
常数、加性与二次部分分别产生 $t^0,t^1,t^2$ 项。
一元多项式对全部有理数为零，故各次项分别为零，不会跨次数抵消。

常数项给出 $sk=0$。加性项展开为
$\sum_jc_jA(a_jx)+\sum_jc_jA(b_jy)$；
分别取 $y=0$、$x=0$ 得到 $\mathcal A_c(V)$ 的两个条件，且它们显然充分。

二次项展开为
\[
 \sum_jc_jB(a_jx,a_jx)
 +2\sum_jc_jB(a_jx,b_jy)
 +\sum_jc_jB(b_jy,b_jy).
\]
先令一个变量为零，对所得对角恒等式极化，得到
\eqref{eq:quadraticconstraints}的前两行。
减去纯项后得到第三行。反过来三行条件使上式为零。
唯一性来自引理\ref{lem:lowdegree}。
\end{proof}

这里的 $K$-线性是指 $\mathcal A_c(V)$、$\mathcal B_c(V)$ 本身的参数空间结构；
它们的成员对输入变量一般只具有 $\Q$-线性或 $\Q$-双线性。

\subsection{两个方向、每个方向恰有两个点}
在四点支撑范围内，没有单点方向恰好意味着非零支撑为两个二点组，
此时不能再出现非零的原点系数 $c_0$。
适当选择两方向上的参数后，方程写成
\begin{equation}\label{eq:twopairs}
 af(x)+bf(tx)+cf(y)+df(uy)=0,
\end{equation}
其中 $a,b,c,d\in K^\times$，$t,u\in K^\times\setminus\{1\}$。
令
\[
 s=a+b+c+d,\qquad \alpha=-a/b,\qquad\beta=-c/d,\qquad
 H=\langle t,u\rangle\le K^\times.
\]

\begin{theorem}[二点组的轨道与字符分类]\label{thm:orbit}
式\eqref{eq:twopairs}的全部解是 $f=k+g$，其中
\[
 sk=0,\qquad g(0)=0,\qquad
 g(tx)=\alpha g(x),\quad g(ux)=\beta g(x).
\]
在 $V\ne0$ 时，非零 $g$ 存在当且仅当
\begin{equation}\label{eq:charactercondition}
 t^m u^n=1\quad\Longrightarrow\quad
 \alpha^m\beta^n=1
 \qquad(m,n\in\Z).
\end{equation}
若该条件成立，它定义唯一的群同态
\[
 \chi:H\to K^\times,\qquad
 \chi(t^mu^n)=\alpha^m\beta^n.
\]
从每个 $H$ 在 $V\setminus\{0\}$ 上的乘法轨道选一个代表 $z$，
任意指定 $g(z)\in K$，并令
\[
 g(hz)=\chi(h)g(z)\qquad(h\in H),
\]
就得到全部解，且这些轨道参数唯一。
若条件失败，则只能有 $g=0$。
\end{theorem}
\begin{proof}
零元组给出 $sf(0)=0$。减去 $k=f(0)$ 后，$g(0)=0$ 且仍满足原方程。
分别令 $y=0$、$x=0$，便得到两个伸缩律；反向相加即得原方程。
由于 $t,u,\alpha,\beta$ 全非零，伸缩律可对任意整数次幂迭代。
若 $g(z)\ne0$，乘法关系 $t^mu^n=1$ 给出
$g(z)=\alpha^m\beta^n g(z)$，所以\eqref{eq:charactercondition}必要。
条件成立时 $\chi$ 定义良好。若 $hz=h'z$ 且 $z\ne0$，
则 $(h-h')z=0$，因而 $h=h'$；乘法作用自由。
故每条轨道上的延拓没有歧义，并给出充分性与完整参数化。
\end{proof}

\begin{example}[确有非广义多项式的齐次解]\label{ex:nongeneral}
取 $K=V=\Q$，考虑
\[
 f(2x)-2f(x)+f(3y)-3f(y)=0.
\]
令 $H=\{2^m3^n:m,n\in\Z\}$，定义
\[
 f(x)=
 \begin{cases}
 x,&x\in H,\\
 0,&x\notin H.
 \end{cases}
\]
则 $f(0)=0$，$f(2x)=2f(x)$、$f(3x)=3f(x)$，所以它确为解。
若 $f$ 是任意有限次数广义多项式，第\ref{sec:models}节的有理一维差分分类表明它是普通多项式。
普通多项式满足 $f(2x)=2f(x)$ 只能含一次项，故必为 $ax$。
但 $f(1)=1$、$f(5)=0$，矛盾。
因此支撑张成二维本身不足以推出广义多项式性。
\end{example}

\subsection{三方向三项方程的半线性分类}
考虑
\begin{equation}\label{eq:threepoint}
 u f(x)+v f(y)+w f(ax+by)+c_0f(0)=0,
\end{equation}
其中 $u,v,w,a,b\ne0$，$c_0$ 可为零。
三个非零方向两两不同。令
\[
 s=u+v+w+c_0,\qquad
 \alpha=-u/w,\qquad \beta=-v/w,\qquad F=\Q(a,b).
\]
由前述分类，解为 $f=k+A$，其中 $sk=0$、$A$ 可加，且
\[
 A(ax)=\alpha A(x),\qquad A(bx)=\beta A(x).
\]

\begin{theorem}[伸缩关系诱导域嵌入]\label{thm:semilinear}
在 $V\ne0$ 时，式\eqref{eq:threepoint}存在非零加性项，当且仅当赋值
$a\mapsto\alpha$、$b\mapsto\beta$ 延拓为域嵌入
$\sigma:F\hookrightarrow K$。
若存在，该嵌入唯一，全部加性项恰为满足
\[
 A(rx)=\sigma(r)A(x)\qquad(r\in F,\ x\in V)
\]
的 $\sigma$-半线性映射。
选取 $V$ 的 $F$-基 $(e_\eta)_{\eta\in E}$，任意指定 $z_\eta\in K$，公式
\[
 A\!\left(\sum_\eta r_\eta e_\eta\right)
 =\sum_\eta\sigma(r_\eta)z_\eta
\]
给出全部解。
\end{theorem}
\begin{proof}
若 $A\ne0$，可加性和两条伸缩律给出
\[
 A(p(a,b)x)=p(\alpha,\beta)A(x)
 \qquad(p\in\Q[X,Y],\ x\in V).
\]
若 $p(a,b)=0$，选 $A(x)\ne0$ 得 $p(\alpha,\beta)=0$。
反过来，若 $p(\alpha,\beta)=0$ 而 $p(a,b)\ne0$，
乘以 $p(a,b)$ 是 $V$ 的双射，上式就迫使 $A=0$，矛盾。
故两个赋值具有完全相同的有理多项式关系。
由此先得到 $\Q[a,b]\to K$ 的单射环同态；
每个非零分母的像非零，所以它唯一延拓到其分式域 $F$。
多项式伸缩律再作用于分母的逆，得到全部 $F$-半线性关系。
反过来，在任意 $F$-基上按所示公式指定值，即得到满足伸缩律的加性映射；
因为 $V\ne0$，至少可以把一个基向量的值指定为 $1$。
\end{proof}


\section{加性与二次成分的张量代数参数化}\label{sec:tensor}
继续使用式\eqref{eq:ranktwo}，定义有限生成域
\[
 F=\Q(a_1,b_1,\ldots,a_m,b_m)\subseteq K.
\]
系数 $c_j$ 属于第一个张量因子 $K$，不必属于 $F$。
下文中的商是交换 $K$-代数商，而商的星号始终是 $K$-线性对偶；
并不要求相应泛函保持乘法。

\subsection{加性成分：一个代数商}
定义
\[
 \mathscr S_1=K\otimes_{\Q}F,\qquad
 u_1=\sum_jc_j\otimes a_j,\qquad
 u_2=\sum_jc_j\otimes b_j,\qquad
 \Uc=\mathscr S_1/(u_1,u_2).
\]
选取 $V$ 的 $F$-基 $(e_\eta)_{\eta\in E}$。

\begin{theorem}[加性项的无冗余参数化]\label{thm:addtensor}
有 $K$-线性同构
\begin{equation}\label{eq:addtensor}
 \mathcal A_c(V)\cong\prod_{\eta\in E}\Uc^*.
\end{equation}
具体地，对任意 $\varphi_\eta\in\Uc^*$，令
\[
 A\!\left(\sum_\eta t_\eta e_\eta\right)
 =\sum_\eta\varphi_\eta\!\left(\overline{1\otimes t_\eta}\right).
\]
这给出全部允许的加性项，参数唯一。
\end{theorem}
\begin{proof}
先固定一个基方向。一个 $\Q$-线性映射 $a:F\to K$ 唯一对应
$K$-线性泛函
\[
 \widetilde a:\mathscr S_1\to K,\qquad
 \widetilde a(k\otimes t)=k\,a(t).
\]
其满足 $\sum_jc_ja(a_jt)=0$ 对所有 $t\in F$，
当且仅当 $\widetilde a(u_1(1\otimes t))=0$ 对所有 $t$。
$\mathscr S_1$ 是所有 $k\otimes t$ 的线性包，故该条件正是
$\widetilde a$ 消灭整个理想 $(u_1)$。
第二条伸缩关系同理，合起来等价于消灭 $(u_1,u_2)$，
即唯一来自 $\Uc$ 上的泛函。
任意加性映射在不同 $F$-基方向上可以独立给定，
而每个向量只涉及有限个基向量，所以所示求和有意义。
反过来限制到每个基方向恢复唯一的 $\varphi_\eta$，得到同构。
\end{proof}

\begin{corollary}\label{cor:addnonzero}
在 $V\ne0$ 时，
\[
 \mathcal A_c(V)=0\quad\Longleftrightarrow\quad\Uc=0.
\]
\end{corollary}
\begin{proof}
若商为零，参数化中所有泛函为零。
若商非零，选择一个非零泛函，并把它指定在至少一个 $F$-基方向上即可。
泛函不可能在全部 $1\otimes t$ 的像上为零，因为这些像张成整个商。
\end{proof}

\subsection{二次成分：带交换对合的代数商}
定义
\[
 \mathscr S_2=K\otimes_{\Q}F\otimes_{\Q}F
\]
以及
\[
 \begin{aligned}
 v_{11}&=\sum_jc_j\otimes a_j\otimes a_j,&
 v_{22}&=\sum_jc_j\otimes b_j\otimes b_j,\\
 v_{12}&=\sum_jc_j\otimes a_j\otimes b_j,&
 v_{21}&=\sum_jc_j\otimes b_j\otimes a_j.
 \end{aligned}
\]
令
\[
 \Wc=\mathscr S_2/(v_{11},v_{22},v_{12},v_{21}).
\]
交换后两个因子的映射固定 $v_{11},v_{22}$，并交换 $v_{12},v_{21}$，
因而在 $\Wc$ 上诱导一个 $K$-线性代数对合 $\iota$。

\begin{theorem}[二次项的无冗余参数化]\label{thm:bitensor}
对每对基指标 $\eta,\zeta$ 选择 $\psi_{\eta\zeta}\in\Wc^*$，
要求
\begin{equation}\label{eq:psisym}
 \psi_{\zeta\eta}=\psi_{\eta\zeta}\circ\iota.
\end{equation}
则公式
\begin{equation}\label{eq:bitensor}
 B\!\left(\sum_\eta t_\eta e_\eta,\sum_\zeta u_\zeta e_\zeta\right)
 =
 \sum_{\eta,\zeta}
 \psi_{\eta\zeta}\!\left(\overline{1\otimes t_\eta\otimes u_\zeta}\right)
\end{equation}
给出全部 $\mathcal B_c(V)$，且参数唯一。
\end{theorem}
\begin{proof}
先固定两个基方向。
$\Q$-双线性映射 $b:F\times F\to K$ 唯一对应
$\mathscr S_2$ 上的 $K$-线性泛函
\[
 \widetilde b(k\otimes t\otimes u)=k\,b(t,u).
\]
例如第一条约束对该基块表示
$\sum_jc_jb(a_jt,a_ju)=0$ 对所有 $t,u$；
它等价于 $\widetilde b$ 消灭 $v_{11}$ 的全部倍数。
其余约束同理给出 $v_{22},v_{12}$。
将混合约束的两输入交换并使用 $B$ 的对称性，又得到
$\sum_jc_jB(b_ju,a_jv)=0$，所以还要消灭 $v_{21}$ 的全部倍数。
因此每个基块恰对应 $\Wc^*$ 的一个元素。
对称性正好等价于\eqref{eq:psisym}，再按两个变量的基展开即得公式。
反向依次验证双线性、对称性和四个理想条件，完成证明。
\end{proof}

\begin{corollary}\label{cor:binonzero}
在 $V\ne0$ 时，
\[
 \mathcal B_c(V)=0\quad\Longleftrightarrow\quad\Wc=0.
\]
\end{corollary}
\begin{proof}
只需证明商非零时存在非零对称型。
因商为含幺代数且非零，$\overline1\ne0$。
取 $\psi\in\Wc^*$ 使 $\psi(\overline1)=1$，并令
$\psi_+=(\psi+\psi\circ\iota)/2$。
则 $\psi_+\circ\iota=\psi_+$，且 $\psi_+(\overline1)=1$。
把它用于一个对角基块，其余基块置零，式\eqref{eq:bitensor}
就给出一个非零对称双线性型。
\end{proof}

\begin{corollary}[三、四方向的单射性]\label{cor:tensorinject}
对 $V\ne0$，
\[
 \begin{aligned}
 r=3:\quad\T_V\text{ 单射}
 &\Longleftrightarrow s\ne0,\ \Uc=0,\\
 r=4:\quad\T_V\text{ 单射}
 &\Longleftrightarrow s\ne0,\ \Uc=0,\ \Wc=0.
 \end{aligned}
\]
\end{corollary}
\begin{proof}
常数、一次广义项和二次广义项的分解唯一。
分别应用定理\ref{thm:exactlowkernel}和前两个非零性判据即可。
对称双线性型非零时其对角函数也非零，因为特征零允许极化。
\end{proof}

这将第\ref{sec:generalideal}节的乘法群代数理想判据，
在当前低支撑范围内转化为两个明确的参数代数商。
它们给出的是全部函数解的结构，并非只描述普通多项式解。

\section{代数参数的有限矩阵与有理矩判据}\label{sec:algebraic}
\subsection{有限维参数代数}
若全部 $a_j,b_j$ 都代数于 $\Q$，则 $F$ 是数域。
记 $d_F=[F:\Q]$，于是
\[
 \dim_K\mathscr S_1=d_F,\qquad
 \dim_K\mathscr S_2=d_F^2.
\]
选取 $F$ 的 $\Q$-基及乘法表，就能构造这些代数的矩阵模型。
对于其中的元素 $z$，以 $M_z$ 表示乘以 $z$ 的 $K$-线性算子。

\begin{proposition}[商维数的秩公式]\label{prop:matrixquotient}
有
\[
 \begin{aligned}
 u:=\dim_K\Uc
 &=d_F-\rank\begin{bmatrix}M_{u_1}&M_{u_2}\end{bmatrix},\\
 w:=\dim_K\Wc
 &=d_F^2-\rank
 \begin{bmatrix}M_{v_{11}}&M_{v_{22}}&M_{v_{12}}&M_{v_{21}}\end{bmatrix}.
 \end{aligned}
\]
\end{proposition}
\begin{proof}
有限维交换代数中，由 $z_1,\ldots,z_k$ 生成的理想，
作为向量空间恰为 $\sum_i\im M_{z_i}$。
它是横向拼接矩阵 $[M_{z_1}\ \cdots\ M_{z_k}]$ 的像。
商的维数等于整个代数维数减去这个像的维数，得到两式。
\end{proof}

令 $w_+=\dim_K\Wc^\iota$。
对合矩阵也可从交换两个 $F$-因子的线性映射直接求出；
由于 $2$ 可逆，$\Wc$ 分解为正、负特征子空间。

\begin{theorem}[有限基方向时的全核维数]\label{thm:dimension}
若另有 $e=\dim_FV<\infty$，则
\[
 \begin{aligned}
 r=3:\qquad
 \dim_K\ker\T_V
 &=\mathbf1_{s=0}+eu,\\
 r=4:\qquad
 \dim_K\ker\T_V
 &=\mathbf1_{s=0}+eu+e w_+ +\binom e2w.
 \end{aligned}
\]
\end{theorem}
\begin{proof}
加性项由 $e$ 个互相独立的 $\Uc^*$ 参数化，维数为 $eu$。
二次项中，每个对角基块要求 $\psi=\psi\circ\iota$，
其参数空间维数为 $w_+$。
每对不同基指标只需任取一个泛函，另一个由对称性决定，故贡献 $w$。
共有 $e$ 个对角块和 $\binom e2$ 个非对角块。
加上可能存在的一维常数空间即得结论。
\end{proof}

若 $e$ 无限，应保留定理\ref{thm:addtensor}、\ref{thm:bitensor}
中的直积参数化，不把这些有限维秩公式解释成无限基数的形式减法。
尤其即使 $V$ 对 $K$ 有限维，$\dim_FV$ 仍可能无限。

\begin{remark}[代数参数与超越参数的区别]
特征零数域扩张可分。取包含 $K$ 的代数闭包 $\Omega$，
以一个本原元的最小多项式分解可得
$\Omega\otimes_{\Q}F\cong\Omega^{d_F}$，
进而 $\Omega\otimes_K\mathscr S_2\cong\Omega^{d_F^2}$。
这些有限乘积代数的理想和商都按坐标分解，商中没有非零幂零元。
由于域扩张保持向量空间嵌入，$\Uc$、$\Wc$ 也没有非零幂零元。
这里用的是域的代数可分性；当参数超越时，商可能含幂零元，
第\ref{sec:derivexample}节将给出明确例子。
\end{remark}

\subsection{有理参数：六个矩决定低支撑全核}
若所有 $a_j,b_j\in\Q$，定义
\[
 M_{pq}=\sum_jc_j a_j^p b_j^q\qquad(p+q>0).
\]
对幂指数为零的因子约定为 $1$。
因为此时 $F=\Q$，有 $\mathscr S_1=\mathscr S_2=K$。

\begin{proposition}[有理参数的完全简化]\label{prop:rationalmoments}
有
\[
 \mathcal A_c(V)=
 \begin{cases}
 \Hom_{\Q}(V,K),&M_{10}=M_{01}=0,\\
 0,&\text{否则},
 \end{cases}
\]
以及
\[
 \mathcal B_c(V)=
 \begin{cases}
 \operatorname{SymBil}_{\Q}(V,K),
   &M_{20}=M_{11}=M_{02}=0,\\
 0,&\text{否则}.
 \end{cases}
\]
因此在三个、四个方向情形，全部函数的齐次核由
$s,M_{10},M_{01},M_{20},M_{11},M_{02}$ 决定，
其中 $r=3$ 时不需要二次矩。
\end{proposition}
\begin{proof}
$\Q$-齐次性给出
$\sum_jc_jA(a_jx)=M_{10}A(x)$，
另一条同理。
双线性条件分别成为
$M_{20}B(u,v)=M_{02}B(u,v)=M_{11}B(u,v)=0$。
域中的任一非零标量都可逆，故任一相应矩非零就迫使整项为零；
矩全零则不再有约束。
\end{proof}

\begin{example}[二次方程的六矩检验]
对
\[
 f(x+y)+f(x-y)-2f(x)-2f(y)=0
\]
有 $s=-2$、$M_{01}=-2$，所以没有常数项或加性项；
而
$M_{20}=M_{11}=M_{02}=0$。
四个支撑方向互异，故全部函数解恰为
$f(x)=B(x,x)$，$B$ 为任意对称 $\Q$-双线性型。
\end{example}

\section{超越参数算例：恒等映射与形式导子}\label{sec:derivexample}
取 $K=V=\Q(\tau)$，其中 $\tau$ 超越于 $\Q$。
考虑完全线性自然的紧线性差
\begin{equation}\label{eq:derivoperator}
 P(f,x,y)=
 \tau^2f(x)+\tau^2f(y)
 -2\tau f(\tau(x+y))+f(\tau^2(x+y)).
\end{equation}
其支撑为 $(1,0),(0,1),(\tau,\tau),(\tau^2,\tau^2)$，
有四个点、三个方向，模板秩为二。

\begin{theorem}[导子算例的全部解]\label{thm:derivkernel}
以 $\partial=d/d\tau$ 表示有理函数域上的形式导子。
则
\[
 \ker\T_K=\{\,x\mapsto ax+b\,\partial x:a,b\in K\,\}.
\]
特别地 $\dim_K\ker\T_K=2$。
\end{theorem}
\begin{proof}
三方向分类表明 $f=k+A$，其中 $A$ 可加。
这里
$s=2\tau^2-2\tau+1\ne0$，故 $k=0$。
利用 $A$ 的可加性，齐次方程等价于
\begin{equation}\label{eq:derivrelation}
 A(\tau^2x)-2\tau A(\tau x)+\tau^2A(x)=0
 \qquad(x\in K).
\end{equation}
定义 $C(x)=A(\tau x)-\tau A(x)$。
式\eqref{eq:derivrelation}给出 $C(\tau x)=\tau C(x)$。
$C$ 同时为 $\Q$-线性，所以与乘以任意 $\Q[\tau]$ 元素的操作交换。
若 $p(\tau)\ne0$，把 $x$ 换成 $x/p(\tau)$ 得
$C(x/p(\tau))=C(x)/p(\tau)$。
因此 $C$ 是 $K$-线性映射，必有 $C(x)=bx$，其中 $b=C(1)$。

形式导子由通常的多项式微分和商法则定义，满足
$\partial(\tau x)-\tau\partial x=x$。
所以 $A_0=A-b\partial$ 与乘以 $\tau$ 交换。
重复前述分母论证，$A_0$ 为 $K$-线性，从而 $A_0(x)=ax$。
反过来 $x\mapsto x$ 和 $\partial$ 均满足\eqref{eq:derivrelation}，
故所有显示函数都是解。
二者在 $1,\tau$ 处的取值矩阵为
$\begin{pmatrix}1&0\\ \tau&1\end{pmatrix}$，行列式为 $1$，故参数唯一。
\end{proof}

这里没有假设未知函数连续或可微。
$\partial$ 作用于输入元素本身的有理函数表达式，是一个确定的域导子。

\subsection{代数商中的二重结构}
把第二个张量因子中的 $\tau$ 记为 $z$。
则 $\mathscr S_1=K\otimes_{\Q}\Q(\tau)$
可视为 $K[z]$ 对所有非零 $\Q[z]$ 多项式作局部化。
两个加性约束的生成元相同，都等于
\[
 z^2-2\tau z+\tau^2=(z-\tau)^2.
\]
因此
\begin{equation}\label{eq:dualnumbers}
 \Uc\cong K[\varepsilon]/(\varepsilon^2),\qquad z\longmapsto\tau+\varepsilon.
\end{equation}
确实，对每个非零 $q\in\Q[z]$，$q(\tau)\ne0$，
所以 $q(\tau+\varepsilon)$ 在该商中为单位，局部化不再改变商。
由形式 Taylor 公式，
\[
 1\otimes r(\tau)\longmapsto r(\tau)+\partial r(\tau)\,\varepsilon
 \qquad(r\in\Q(\tau)).
\]
两个坐标泛函分别对应恒等映射和形式导子。
这说明不能仅用域嵌入描述所有超越参数的加性解；
商中的非零幂零元记录了导数项。
域自同构、微分算子与线性函数方程的关系可参见\cite{KissLaczkovich}；
本算例的完整性已由上面的直接证明给出。

\subsection{非齐次自由度与两个初值}
若 $\T_Kf_0=h$，则所有解为
\[
 f=f_0+ax+b\,\partial x.
\]
若指定 $f(1)=p$、$f(\tau)=q$，唯一参数为
\[
 a=p-f_0(1),\qquad
 b=q-f_0(\tau)-\tau\bigl(p-f_0(1)\bigr).
\]
因此任意两个这样的初值都能在一个非空解仿射空间中选出唯一解。
这不表示任意 $h$ 都可解；$h$ 本身仍须满足有限关系判据。

\section{代数参数算例：共轭嵌入与范数}\label{sec:normexample}
取 $K=V=\Q(\alpha)$，$\alpha^2=2$，并记非平凡共轭为
$\sigma(\alpha)=-\alpha$。
考虑
\begin{equation}\label{eq:normoperator}
 \begin{split}
 P(f,x,y)={}&f(x+(1+\alpha)y)-f(x+y)\\
 &-f(x+\alpha y)+f(x).
 \end{split}
\end{equation}
四个支撑 $(1,1+\alpha),(1,1),(1,\alpha),(1,0)$
方向两两不同，故全部齐次解至多为二次广义多项式。

\begin{theorem}[共轭与范数算例的全部解]\label{thm:normkernel}
有
\[
 \ker\T_K
 =\{\,x\mapsto k+ax+b\sigma(x)+c\,x\sigma(x):
 k,a,b,c\in K\,\},
\]
且四个参数唯一，$\dim_K\ker\T_K=4$。
\end{theorem}
\begin{proof}
按定理\ref{thm:exactlowkernel}写成 $f=k+A+B(x,x)$。
系数总和为零，故常数项自由。
对任意加性 $A$，用
$A((1+\alpha)y)=A(y)+A(\alpha y)$
即可直接验证它满足齐次方程。

对二次部分逐项展开，所有含 $x$ 的项抵消，剩余
$2B(y,\alpha y)$。
因此其充要条件是
\[
 B(y,\alpha y)=0\qquad(y\in K).
\]
在 $\Q$-基 $1,\alpha$ 下置
$B(1,1)=u$、$B(1,\alpha)=v$、$B(\alpha,\alpha)=w$。
对 $y=r+t\alpha$，$r,t\in\Q$，
\[
 B(y,\alpha y)=v r^2+(2u+w)rt+2v t^2.
\]
恒为零当且仅当 $v=0$、$w=-2u$。
故 $B(y,y)=u(r^2-2t^2)=u\,y\sigma(y)$。

任意 $\Q$-线性映射 $A:K\to K$ 由 $A(1),A(\alpha)$ 决定，
唯一写成 $ax+b\sigma(x)$，其中
\[
 a=\frac12\left(A(1)+\frac{A(\alpha)}{\alpha}\right),\qquad
 b=\frac12\left(A(1)-\frac{A(\alpha)}{\alpha}\right).
\]
由广义次数分解的唯一性，常数、加性与二次参数互不混淆。
反向代入已经包含在上述计算中，故得到全部解。
\end{proof}

范数项 $x\sigma(x)$ 不是普通的 $K$-二次单项式 $cx^2$：
若二者恒等，由 $x=1$ 得 $c=1$，而 $x=\alpha$ 给出 $-2=2$，矛盾。
它是底层 $\Q$-结构上的二次广义多项式。

\subsection{张量代数的交叉检验}
这里 $F=K$、$e=\dim_FV=1$。
$\mathscr S_1$ 按两个嵌入 $\id,\sigma$ 分解为 $K\times K$。
两个加性生成元均为零，所以 $\Uc\cong K^2$，$u=2$。

$\mathscr S_2$ 按嵌入对 $(\rho,\theta)$ 分解为四个 $K$-因子。
$v_{11},v_{12},v_{21}$ 为零，而 $v_{22}$ 在 $(\rho,\theta)$ 处的值是
\[
 (1+\rho(\alpha))(1+\theta(\alpha))
 -1-\rho(\alpha)\theta(\alpha)
 =\rho(\alpha)+\theta(\alpha).
\]
同嵌入的两个坐标分别为 $2\alpha,-2\alpha$，被理想消去；
异嵌入的两个坐标为零，保留在商中。
于是 $\Wc\cong K^2$，对合交换两个坐标，故 $w=2$、$w_+=1$。
定理\ref{thm:dimension}给出
$1+2+1=4$，与直接求解完全一致。

\subsection{四个初值与自然性的边界}
对基函数 $1,x,\sigma(x),x\sigma(x)$，
在 $0,1,-1,\alpha$ 四点的取值矩阵为
\[
 \begin{pmatrix}
 1&0&0&0\\
 1&1&1&1\\
 1&-1&-1&1\\
 1&\alpha&-\alpha&-2
 \end{pmatrix},
 \qquad\det=4\alpha\ne0.
\]
故若存在特解 $f_0$，任意指定这四点的函数值都可唯一消除全部自由参数。

式\eqref{eq:normoperator}属于完全线性自然模板；
但四个支点的坐标和分别为 $2+\alpha,2,1+\alpha,1$，
并非全部为 $1$。
由推论\ref{cor:affine}，该族不具有对两个参数同时平移的仿射自然性。
这两种自然性要求不能混用。

\section{回到非齐次问题：存在性、全部解与初值}\label{sec:inhomfinal}
\subsection{秩二模板的精确相容判据}
对式\eqref{eq:ranktwo}，在自由向量空间 $D_V=K^{(V)}$ 中定义
\[
 r_{x,y}=c_0[0]+\sum_{j=1}^m c_j[a_jx+b_jy],\qquad
 L_c(V)=\Span_K\{r_{x,y}:x,y\in V\}.
\]
若某些实际取值点重合，必须先在该行中合并它们的系数；
这是行向量的定义本身，并非额外假设。

\begin{theorem}[本报告范围内的非齐次解结构]\label{thm:inhomfinal}
给定任意 $h:V^2\to K$，方程 $P(f,x,y)=h(x,y)$ 有解，
当且仅当每条有限行关系都由 $h$ 保持，即
\[
 \sum_{\ell=1}^N d_\ell r_{x_\ell,y_\ell}=0
 \quad\Longrightarrow\quad
 \sum_{\ell=1}^N d_\ell h(x_\ell,y_\ell)=0.
\]
若有一个解 $f_0$，全部解为 $f_0+\ker\T_V$，
其中 $\ker\T_V$ 恰按第\ref{sec:kernelclassification}节的四类分类给出。
不相容时必有一条有限关系作为不可解证书。
\end{theorem}
\begin{proof}
相容性允许定义
$\varphi_h(\sum d_\ell r_{x_\ell,y_\ell})
=\sum d_\ell h(x_\ell,y_\ell)$。
将这个 $L_c(V)$ 上的泛函延拓到 $D_V$，
并令 $f(x)$ 为延拓泛函在 $[x]$ 上的值，即得特解。
反向把任一解作用到有限行关系上就得到相容性。
全部解是一个特解加齐次核，核的分类已在前文逐一证明。
若条件不成立，其失败按定义就由某条有限关系见证。
\end{proof}

该定理是对任意右端的完整存在性判据，
但并不声称一般模板的全部行关系已化为有限张列表。
两个参数代数商 $\Uc,\Wc$ 直接控制的是齐次自由度；
非齐次存在性仍由上述行关系决定。
若模板由原来的 $n$ 个参数约化而来，还须先检查
第\ref{sec:lowrank}节的右端纤维条件。

\subsection{有限维齐次核总可以用同样多的点值消除}
\begin{proposition}[最少初值数]\label{prop:initialvalues}
设某一标量值线性函数方程的齐次核 $Z\subseteq K^V$
满足 $\dim_KZ=d<\infty$。
则存在互异的点 $x_1,\ldots,x_d\in V$，使
\[
 Z\longrightarrow K^d,\qquad
 g\longmapsto(g(x_1),\ldots,g(x_d))
\]
为同构。
因此对每个可解右端，任意给定这 $d$ 个点值都选出唯一解。
少于 $d$ 个点值不能消除全部齐次自由度。
\end{proposition}
\begin{proof}
所有点值泛函在 $Z$ 上分离向量：若 $g(x)=0$ 对每个 $x$ 成立，则 $g=0$。
若它们的线性包不是 $Z^*$，有限维对偶性将给出一个非零 $g$
被所有这些泛函消灭，矛盾。
因此可从点值泛函中选取 $Z^*$ 的一组基；
相应的 $d$ 个点就是所求点，取值映射为同构。
在解仿射空间 $f_0+Z$ 上，指定值只是在这个同构中指定一个向量，
所以任意值都能且只能实现一次。
从 $d$ 维空间到 $K^m$ 的线性映射在 $m<d$ 时不可能单射，给出最少性。
\end{proof}

\subsection{退化空间}
若 $V=0$、$W=K$ 且 $n\ge1$，定义域及参数域都只有一个点。
任何标量自然模板都退化为
\[
 s f(0)=h(0,\ldots,0).
\]
$s\ne0$ 时对每个右端唯一可解；
$s=0$ 时仅零右端可解，且 $f(0)$ 任意。
所以定理\ref{thm:orbit}、\ref{thm:semilinear}
关于非零解的存在条件，以及 $\Uc,\Wc$ 的非零检测，
均明确要求 $V\ne0$。
若值域本身为零空间，则只有一个函数及一个右端，解总唯一。

\section{结论边界、核查说明与后续问题}\label{sec:scope}
\subsection{已经获得的完整结论}
本报告的核心结论可归纳为如下三层，每层的假设均已在正文给出。
\begin{enumerate}
\item 对紧标量线性差，全部非齐次存在性障碍恰为有限行关系；
核为行空间余核的完整代数对偶。
有限维值域可作相同的对偶处理；一般无限维算子系数须保留延拓障碍。
\item 对全部线性映射自然的紧算子族，通用空间 $K^n$
给出唯一的有限支集模板。
一般单射性可以在一维上检测，并由乘法群代数理想刻画；
一元模板、有限域及普通多项式类各有更明确的模型。
\item 在特征零、模板秩二、总支撑至多四点时，
所有齐次函数解已按方向分布完整分类。
单点方向控制广义次数；三个、四个方向的允许成分分别由
$\Uc,\Wc$ 参数化；没有单点方向的唯一构型由轨道字符控制。
\end{enumerate}

\subsection{本次整理明确保留的适用限制}
\begin{enumerate}
\item 线性差是对函数变量的 $K$-线性；这不强迫未知函数对 $V$ 的输入为 $K$-线性。
“加性”在特征零下只自动给出 $\Q$-线性。
\item 支撑点数包含原点，方向数不包含原点。
模板秩二及总支撑四点是完整低支撑分类的边界，
不应把其结论无条件推广到任意支撑。
\item 三个、四个方向产生的是广义多项式，而非普通 $K$-多项式。
有理参数使矩判据控制全部加性、双加性成分；
一般参数可能引入共轭和导子。
\item 参数代数商的对偶是线性泛函空间。
用环同态或域嵌入替代它，会遗漏导子算例中的非约化成分。
\item 只有在 $F/\Q$ 有限且 $\dim_FV$ 有限时，
前述公式同时给出有限的全核维数。
单射性判据本身不要求 $\dim_FV$ 有限。
\item 核的完整分类与任意右端可解不是同一个结论。
初值唯一化都以给定右端已经可解为前提。
\end{enumerate}

\subsection{可继续研究的方向}
当前四点分类已经封闭；更高支撑需要新增结构。
自然的后续问题包括：多个重复方向同时存在时，
方向合并函数与原函数之间如何耦合；
高阶广义多项式能否统一用带对称群作用的多重张量代数商参数化；
在特定模板下能否将非齐次行关系化为简洁的余循环恒等式；
以及超越参数代数商的局部结构如何控制更高阶形式微分项。
这些问题未被本报告的四点结论自动解决。

\medskip
\noindent\textbf{文献使用说明。}
本报告的主要论证均在正文完成。
下面列出的文献只用于标识线性函数方程、广义多项式与谱综合的相关研究背景；
本报告不以未在正文证明的一般谱综合结论作为低支撑分类的前提。
广义指数多项式及其与可分解函数的关系可参见\cite{Laczkovich2018}，
但本报告所需的零至二次分解使用引理\ref{lem:lowdegree}的直接证明。

\begin{thebibliography}{9}
\raggedright
\addcontentsline{toc}{section}{参考文献}
\bibitem{KissVarga}
G. Kiss and A. Varga.
\textit{Existence of nontrivial solutions of linear functional equations}.
Aequationes Mathematicae \textbf{88} (2014), 151--162.
\url{https://doi.org/10.1007/s00010-013-0212-z}.

\bibitem{KissLaczkovich}
G. Kiss and M. Laczkovich.
\textit{Linear functional equations, differential operators and spectral synthesis}.
Aequationes Mathematicae \textbf{89} (2015), 301--328.
\url{https://doi.org/10.1007/s00010-014-0269-3}.

\bibitem{Laczkovich2018}
M. Laczkovich.
\textit{A characterization of generalized exponential polynomials in terms of decomposable functions}.
arXiv:1812.06434 (2018).
\url{https://arxiv.org/abs/1812.06434}.
\end{thebibliography}

\end{document}

